\documentclass[11pt,letterpaper]{article}
\usepackage[T1]{fontenc}
\usepackage[utf8]{inputenc}
\usepackage{lmodern}
\usepackage[margin=1in,headheight=15pt]{geometry}
\usepackage{microtype}
\usepackage{amsmath,amssymb,amsthm,mathtools}
\usepackage{booktabs,longtable,array,enumitem,multicol}
\usepackage[dvipsnames]{xcolor}
\usepackage{fancyhdr}
\usepackage{hyperref}
\usepackage{listings}
\definecolor{inkblue}{HTML}{17365D}
\definecolor{muted}{HTML}{536471}
\hypersetup{colorlinks=true,linkcolor=inkblue,citecolor=inkblue,urlcolor=inkblue,
 pdftitle={Computable Surreal Numbers: A Reconciled Theory of Structural and Numerical Computability},
 pdfauthor={Unified synthesis of three supplied reports}}
\urlstyle{same}
\pagestyle{fancy}
\fancyhf{}
\fancyhead[L]{\small\textsc{Computable surreal numbers}}
\fancyhead[R]{\small\textsc{Reconciled research report}}
\fancyfoot[C]{\thepage}
\renewcommand{\headrulewidth}{0.3pt}
\setlength{\parskip}{4pt}
\setlength{\parindent}{1.2em}
\setlength{\emergencystretch}{2em}
\setcounter{tocdepth}{1}
\numberwithin{equation}{section}
\newtheorem{theorem}{Theorem}[section]
\newtheorem{proposition}[theorem]{Proposition}
\newtheorem{lemma}[theorem]{Lemma}
\newtheorem{corollary}[theorem]{Corollary}
\theoremstyle{definition}
\newtheorem{definition}[theorem]{Definition}
\newtheorem{example}[theorem]{Example}
\theoremstyle{remark}
\newtheorem{remark}[theorem]{Remark}
\newcommand{\PP}{\mathcal P}
\newcommand{\No}{\mathbf{No}}
\newcommand{\Rc}{\mathbb R_{\mathrm c}}
\newcommand{\Cc}{\mathbb C_{\mathrm c}}
\newcommand{\Sc}{\mathcal S_{\mathrm{str}}}
\newcommand{\Lc}{\mathcal L_{\mathrm c}}
\newcommand{\Pc}{\mathcal P_{\mathrm c}}
\newcommand{\Q}{\mathbb Q}
\newcommand{\R}{\mathbb R}
\newcommand{\N}{\mathbb N}
\newcommand{\Z}{\mathbb Z}
\newcommand{\supp}{\operatorname{supp}}
\newcommand{\st}{\operatorname{st}}
\newcommand{\res}{\operatorname{res}}
\newcommand{\val}{\operatorname{v}}
\newcommand{\coef}[2]{[t^{#1}]#2}
\newcommand{\cut}[2]{\left\{#1\,\middle|\,#2\right\}}
\newcommand{\CK}{\omega_1^{\mathrm{CK}}}
\newcommand{\Oh}{\mathcal O}
\newcommand{\mh}{\mathfrak m}
\DeclareMathOperator{\lcm}{lcm}
\lstset{basicstyle=\small\ttfamily,columns=fullflexible,breaklines=true,
 frame=single,rulecolor=\color{muted},showstringspaces=false}

\begin{document}
\hypersetup{pageanchor=false}%
\begin{titlepage}
\vspace*{0.8in}
\noindent{\color{inkblue}\Large\textsc{Foundations and computation}}
\vspace{0.35in}

\noindent{\Huge\bfseries Computable\\[5pt] Surreal Numbers}
\vspace{0.25in}

\noindent{\LARGE Structural names, Puiseux series,\\[4pt] and effective left-finite fields}
\vspace{0.35in}

\noindent{\large A unified mathematical report with proofs,\\
representation contracts, and a reconciliation ledger}
\vspace{0.45in}

\noindent\textbf{Merged and reconciled from three supplied reports}\\
Prepared for the Surreal project \;\textbullet\; 21 September 2026

\vfill
\noindent\begin{minipage}{\textwidth}
\textbf{Central distinction.}
A program describing a well-founded surreal construction need not compute
its real value in the ordinary numerical sense. Conversely, a coefficient
algorithm need not reveal the finite support information needed for multiplication.
This article develops both notions without identifying them.

\medskip
\textbf{Status.} Common material is consolidated; distinct representations and
complementary proofs are retained. The reconciliation ledger records source
locations, strengthened formulations, and corrections. Established results
remain attributed to prior work; no new priority claim is made.
The article has not been refereed or formalized in Lean. Its accompanying
finite exact-arithmetic checks are tests, not machine-checked proofs.
\end{minipage}
\end{titlepage}
\hypersetup{pageanchor=true}%

\begin{abstract}
This report merges three accounts of computable surreal numbers into one
representation-sensitive theory. Hereditarily effective Conway-cut names define
the established structural field $\Sc$, whose real trace is the hyperarithmetic
reals. For ordinary numerical access, two different fields must be retained:
$\Pc$, using a lower-bounded lattice with one ramification denominator per
name, and $\Lc$, using rational left-finite supports equipped with computable
finite candidate covers. Both are countable real-closed subfields of $\No$
with real trace exactly $\Rc$, and $\Pc\subsetneq\Lc$. The larger field is the
effective sequential completion of the smaller for a supplied exact valuation
stabilization modulus; neither assertion is an unqualified classical
completeness claim.

The unified proofs give uniform ring operations, exact symbolic truncation,
certified reciprocal and root algorithms, effective Hensel lifting, formal
functional calculus, and finite-rank extensions. Two independent routes to
Puiseux real closedness are preserved: a henselian valuation argument and a
Newton--Puiseux argument with an explicit finite-jet lifting lemma. The latter
explains why finitely many branch data suffice for a single coefficient
algorithm, without assuming a uniform branch selector.

The obstructions are equally explicit. Computable Hahn coefficients can have
a noncomputable product, even with decidable left-finite supports; a separate
example has every convolution fiber of size at most one. One halting oracle
is sharp for unrestricted coefficient convolution and for several normalization
tasks, but its output role differs in the two settings. Equality, sign,
inversion, roots, and modulus are not silently conflated with extensional
closure. Oracle-relative fields recover the Turing-degree ordering. A source
crosswalk, a corrected nonzero-polynomial hypothesis, repository-specific
implementation contracts, and independently rerun finite checks complete the
reconciliation. None of the finite tests is a formal proof of the infinite
mathematical statements.
\end{abstract}
\begingroup
\small
\setlength{\parskip}{0pt}
\begin{multicols}{2}
\tableofcontents
\end{multicols}
\endgroup
\newpage

\section{The question and the reconciled answer}
\label{sec:introduction}

A computable real number is not a real number whose entire decimal expansion
has already been stored. It is one for which a finite algorithm can supply
arbitrarily accurate rational approximations with an effective error bound.
Replacing $\R$ by $\No$ changes much more than the possible magnitude of the
answer. A surreal has an ordinal-length sign sequence and a transfinite
normal form. There is no single sequence of positive surreal error tolerances
cofinal toward zero in the whole class. Moreover, the algebra of infinite
normal forms depends on support information that individual coefficient
queries need not provide.

The appropriate conclusion is not that computation is impossible, but that
\emph{the representation is part of the mathematical definition}. We first separate structural from numerical effectivity:
\[
 \Sc=\text{structurally computable surreals},\qquad
 \Lc=\text{numerically computable left-finite surreals}.
\]
The first is broad and hereditary. The second is a precisely specified
rank-one numerical workspace containing the distinct Puiseux model $\Pc$.
None is advertised as the unique possible meaning of ``computable surreal.''


\subsection{Source identities and the central reconciliation}
Throughout, \textbf{A} denotes the report in
\path{computable_surreal_numbers.zip}; \textbf{B} the report in
\path{computable_surreal_numbers(1).zip}; and \textbf{C} the report in
\path{computable_surreal_numbers(2).zip}. The suffixes are upload identifiers,
not chronological or mathematical version numbers. Their source manuscripts
are \emph{not} redistributed with this report; the accompanying provenance
manifest records their checksums and line counts.
Appendix~\ref{app:reconcile} provides the topic and correction crosswalk.

A uses a larger effective left-finite field as its main numerical model;
B and C use a bounded-denominator Puiseux field. They therefore do not make
inconsistent assertions about one field. We keep both symbols and prove the
inclusion explicitly. In particular, C's request for an effective support
enlargement beyond Puiseux series is answered, at the left-finite level, by
A's finite-cover construction. The question left open here concerns more
general supports, not that already constructed enlargement.

The strongest compatible statements have been selected rather than averaged:
C's full validity classification and degree-preserving sign example; A's
left-finite support counterexample, Hensel construction, and effective
completion; B's sharp convolution oracle bound, algebraic-branch obstruction,
and differential consequences; and B--C's finite-jet argument and finite-rank
Puiseux construction. C's oracle hierarchy is extended below to both numerical
models by the same explicit proof. These are synthesis choices, not assertions
of independent novelty.

\begin{center}
\renewcommand{\arraystretch}{1.17}
\begin{tabular}{@{}p{.16\textwidth}p{.29\textwidth}p{.22\textwidth}p{.24\textwidth}@{}}
\toprule
Model & Information supplied by a name & Real trace & Main qualification\\
\midrule
$\Sc$ & Hereditary enumerations of separated cuts & Hyperarithmetic reals &
Semantic validity is $\Pi^1_1$-complete.\\
$\Pc$ & Uniform real coefficients on one lower-bounded rational lattice &
$\Rc$ & Computable-modulus limits may escape the lattice.\\
$\Lc$ & Uniform real coefficients and complete finite support covers &
$\Rc$ & Effective completeness requires a supplied modulus.\\
\bottomrule
\end{tabular}
\end{center}
Both numerical models have uniform ring operations but not a uniform raw-name
inverse. Neither real trace nor algebraic closedness supplies a missing
representation converter. We do not assert a general computable translation
between structural and numerical names.

\subsection{Prior work and what is being developed}

Recursive surreal numbers were studied by Harkleroad in 1990 and several
notions were compared by Lurie in 1998 \cite{harkleroad,lurie}. Later
presentations by Hamkins, reporting joint work with Turetsky, explicitly
identify the hereditary structural version with hyperarithmetic sign
sequences and establish its real closedness \cite{hamkins,hamkins2025,turetsky}.
Those results are \emph{prior theory}, not discoveries of this article.

The numerical carrier used below is equally classical: real series with
rational, left-finite support form the Levi--Civita field. A recent primary
source recalls this definition and its classical algebraic properties
\cite[\S1.1]{restrepo}. Our question is which of its elements and operations
are effective under a representation that continues to accept ordinary
Cauchy names for real constants. The effective finite-cover requirement,
the explicit counterexamples, and the resulting collection of closure and
nonuniformity proofs are the focus here. We do not infer novelty merely
from having supplied a proof; related effective valued-field formulations
may exist under other representations or terminology.

The standard background is Conway's construction and Gonshor's normal-form
theory \cite{conway,gonshor}. Only one substantial valuation-theoretic input
is used in the numerical real-closedness proof: a henselian finite extension
in residue characteristic zero has degree equal to ramification index times
residue degree. We isolate that input explicitly in
Section~\ref{sec:realclosed}.


The independent finite-jet proof in Section~\ref{C-sec:realclosed} instead uses
the characteristic-zero Newton--Puiseux theorem, stated explicitly in the
introductions of \cite{rond,kedlaya}. Its additional effective step is proved,
not inferred from the computability of each coefficient separately. The
user-supplied overview \cite{wiki} motivates the ordinary-real baseline;
the report fixes fast Cauchy names rather than treating decimal, Dedekind-cut,
and Cauchy encodings as uniformly interchangeable. Objectwise equivalence
of computable values is weaker than a uniform conversion of names.

All algorithms here are ordinary finite-time Turing computations per query.
The ordinal length of a support is not a machine runtime. Infinitary-machine
approaches, such as the $\kappa$-time and $\kappa$-tape framework of
Galeotti--Nobrega \cite{galeotti}, change the computational model and are
not counterexamples to the ordinary-computability obstructions below.

\subsection{A map of the main results}

The main numerical theorem, assembled from the proofs below, is as follows.

\begin{theorem}[Numerically conservative effective core]
\label{thm:main}
Let $\Lc$ be the class of surreal normal forms with an effective left-finite
Cauchy--Hahn name, as in Definition~\ref{def:elf}. Then:
\begin{enumerate}[label=\textup{(\roman*)},leftmargin=2em]
\item $\Lc$ is a countable real-closed subfield of $\No$, and
$\Lc\cap\R=\Rc$. Its natural valuation has value group $\Q$ and residue
field $\Rc$.
\item Addition, subtraction, multiplication, rational monomial shifts,
and coefficient extraction are computable uniformly on names.
\item Every nonzero element has an inverse in $\Lc$. Computing an inverse
uniformly is possible when the leading exponent is supplied, and is
impossible on all nonzero raw names without additional information.
\item Simple residue roots lift effectively. The computable Puiseux field
$\Pc$ is a proper real-closed subfield of $\Lc$, and $\Lc$ is its effective
sequential completion for exact valuation stabilization.
\item $\Lc[i]$ is algebraically closed. Formal evaluation at infinitesimals,
formal differentiation, and residue-free formal integration admit explicit
algorithms.
\end{enumerate}
\end{theorem}

This theorem does \emph{not} say that $\Lc$, with its raw representation, is
a computable ordered field in the sense of having decidable equality and
uniform field operations. Indeed, part (iii) rules out that interpretation.
The distinction between an extensional algebraic closure property and a
uniform operation on names will be maintained throughout.

Two negative results explain why the qualifications in this theorem are
substantive. Theorem~\ref{thm:badproduct} constructs computable coefficient
functions, even with decidable supports, whose Hahn product has a
noncomputable coefficient function. Theorem~\ref{thm:nogo} shows that
uniform inversion already conflicts with ordinary real input and coefficient
access on the rational-function family $a+t$.

\section{Conventions, names, and three topologies}
\label{sec:conventions}

\subsection{Surreal and Hahn conventions}

We work classically, for example in ZFC with definable classes, or in a
universe arrangement in which the full surreal carrier is one universe
larger than the small option families. Every individual support and every
program considered here is a set. We never quantify over a purported set of
all surreals.

For small separated sets $L,R\subseteq\No$, $\cut LR$ denotes the
\emph{simplest} surreal strictly between every member of $L$ and every
member of $R$. It is not an ordinary supremum. The birthday $b(x)$ is the
length of the canonical sign sequence of $x$.

Fix
\[
 t=\omega^{-1},\qquad t^q=\omega^{-q}\quad(q\in\Q).
\]
Under Conway normal forms, the usual Hahn field $\R((t^\Q))$, with
well-ordered support in the increasing exponent convention, is an ordered
subfield of $\No$. All uses of this identification are classical consequences
of normal-form theory \cite{conway,gonshor}; we do not claim that the
repository already formalizes this bridge. A series is written
\[
 x=\sum_{q\in\Q}a_qt^q,\qquad
 \val(x)=\min\supp(x)\quad(x\ne0),\qquad \val(0)=+\infty.
\]
Its sign is the sign of $a_{\val(x)}$. Larger valuation means a smaller
non-Archimedean scale. We set $\coef qx=a_q$.

A set $S\subseteq\Q$ is \emph{left-finite} if $S\cap(-\infty,B)$ is
finite for every rational $B$. Every nonempty left-finite set is bounded
below and has a least element: choose one member $s$, and inspect the finite
nonempty set $S\cap(-\infty,s+1)$. The same argument applies to each nonempty
subset of $S$, so $S$ is well ordered. An infinite left-finite support has
order type $\omega$ and tends to $+\infty$.

\subsection{What a computable real coefficient means}

Write $\Rc$ for the ordinary computable reals. Our real names are functions
$A:\N\to\Q$ satisfying
\begin{equation}\label{eq:cauchyname}
 |A(k)-a|\le 2^{-k}.
\end{equation}
For a family $a_q$ we require a \emph{single} algorithm $A(q,k)$ satisfying
this bound jointly in the rational index $q$ and precision $k$. Saying
merely that each $a_q$ is a computable real is weaker: the family
$a_n=\chi_K(n)$ consists of rational numbers but is not uniformly computable.
Here $K$ is a fixed undecidable computably enumerable halting set.

A name is a finite program plus any explicitly specified finite data.
Correctness conditions on a name are semantic promises. There is no
assumption that an interpreter can verify totality, all error inequalities,
or all support conditions. A \emph{uniform operation} is an algorithm which
transforms every valid input name satisfying its stated promises into a valid
output name. An \emph{extensional closure theorem} only asserts that every
resulting value has some computable name.

\subsection{Why full-surreal Cauchy approximation is the wrong starting point}

\begin{proposition}[No small nontrivial convergence in the fine order topology]
\label{prop:topology}
A set-indexed net in $\No$ which converges in the full order topology is
eventually equal to its limit. A set-indexed Cauchy net for its ordered-field
uniformity is eventually constant.
\end{proposition}
\begin{proof}
For a net $(x_i)_{i\in I}$ and a proposed limit $x$, form the set
$D=\{|x_i-x|:i\in I,\ x_i\ne x\}$ of positive distances. The small cut
$\varepsilon=\cut 0D$ is positive and smaller than every member of $D$.
Convergence eventually gives $|x_i-x|<\varepsilon$, forcing $x_i=x$.
For the Cauchy assertion use instead the set of all nonzero pairwise
distances. The Cauchy condition makes all sufficiently late pairwise
differences zero. Empty distance sets cause no difficulty.
\end{proof}

This is consistent with the topology results documented in the repository
\cite{repo}. It does not prevent computation in a fixed workspace. Inside
the Levi--Civita field, $t^n$ tends to zero for its \emph{intrinsic}
valuation topology, since $n$ eventually exceeds every rational valuation.
In $\No$, however, $t^n>\omega^{-\omega}$ for every integer $n$.
Indeed, the subspace topology that the full surreal order induces on the
rational-exponent field is discrete, whereas its intrinsic order topology
is not.

There are thus three distinct notions in play. Real coefficient errors are
measured by the ordinary real absolute value. Valuation truncation preserves
coefficients \emph{exactly} below an exponent bound. Full-surreal order
neighborhoods allow all surreal scales. Our numerical names combine the
first two; they do not make the full surreal topology an effective metric
space. Nor are they an admissible presentation of the pure valuation metric
with the usual Cauchy representation on constants: nonzero real constants
all have valuation zero, however small their real magnitudes.

\section{The established structural notion}
\label{sec:structural}

\subsection{Hereditary effective cuts}

\begin{definition}[Structural presentation]
\label{def:structural}
A raw structural name is a program index $e$ whose execution enumerates two
collections of program indices, called its left and right options. Unfold
these indices hereditarily into a rooted option tree. The name is
\emph{valid} if this tree is well founded and, at every node, the values of
all left options are strictly below the values of all right options.
Well-founded recursion assigns to the node the simplest surreal separating
its options. The set $\Sc$ consists of the values of valid names.
\end{definition}

An enumeration may never announce that it has finished. A leaf can be a
program that runs forever without output. Thus validity is not the same as
termination of the enumerating program. Repetition of an option has no
semantic effect. A finite program describes a well-founded infinite object;
no ordinary Turing machine is asserted to execute one step at every ordinal
stage of its birthday.

The following result is an external classification, included to locate our
work accurately rather than reproved here.

\begin{theorem}[Known structural classification]
\label{thm:known}
For hereditary effective surreal presentations, the values are exactly the
surreals with hyperarithmetic sign sequences, interpreted on effective
ordinal presentations. Equivalently, with the usual sign-sequence coding,
they are $\No\cap L_{\CK}$. This is a real-closed field. Its real slice is
the hyperarithmetic reals, not merely $\Rc$.
\end{theorem}

The effectivization program originates in
\cite{harkleroad,lurie}; the displayed modern characterization and
real-closedness statement are given in \cite{hamkins,hamkins2025,turetsky}.
The full Lurie article was not accessible during preparation; its published
abstract and metadata were checked, while the explicit characterization
was checked against the later authors' public statements and lecture notes.
Here hyperarithmetic means computable by effective transfinite recursion
below the Church--Kleene ordinal; it is strictly broader than ordinary
Turing computability. We use Theorem~\ref{thm:known} only for the structural
comparison, not in any proof of the numerical main theorem.

\subsection{Elementary closure mechanisms}

\begin{proposition}[Structural ring compilation and computable cuts]
\label{prop:structring}
Negation, addition, and multiplication admit uniform compilers on valid
structural names. If $(x_n)$ and $(y_n)$ are uniformly named families with
$x_n<y_m$ for all $n,m$, then $\cut{\{x_n:n\in\N\}}{\{y_n:n\in\N\}}$
has a uniformly constructed valid structural name.
\end{proposition}
\begin{proof}
For cuts, produce one root program which dovetails the two enumerations of
names. Every branch below the root lies in a valid input tree; hence the
new tree is well founded, and the promised inequalities give separation.

For arithmetic use the Conway option formulas. Negation interchanges sides
and negates each option. Addition emits $x^L+y$ and $x+y^L$ on the left,
and $x^R+y$ and $x+y^R$ on the right. Multiplication emits
\[
 x^Ly+xy^L-x^Ly^L,\quad x^Ry+xy^R-x^Ry^R
\]
on the left, and the corresponding mixed left--right expressions on the
right. These are finite syntactic expressions in recursively named smaller
products. Standard effective self-reference supplies the indices of the
recursive compilers. Well-founded induction on the input option trees,
using the well-founded product order for pairs of operands, proves that the
resulting finite-expression unfoldings are valid and have the intended
values. Separation is exactly the classical Conway arithmetic theorem
\cite{conway}. The construction never decides which of two arbitrary input
values is larger.
\end{proof}

This proof does not give a blind reciprocal compiler. The known assertion
that $\Sc$ is a field must not be substituted for such an algorithm. In
particular, restrictions to positive options in classical reciprocal
formulas cannot simply be implemented by assuming a decidable sign test.

\begin{proposition}[Computable ordinals and lack of an exhaustive enumeration]
\label{prop:structcount}
Every computable ordinal belongs to $\Sc$. The set $\Sc$ is countable,
but no uniform sequence of valid structural names is cofinal in it. In
particular, no such sequence enumerates all its values.
\end{proposition}
\begin{proof}
Given a computable well-order $(D,\prec)$, assign to each $d\in D$ the
all-left name whose options are the names for $c\prec d$, and place all
these names below a new all-left root. Induction identifies the root's value
with the order type of $D$. The compilers are effective since $D$ and
$\prec$ are decidable.

Countability follows from countably many program indices. If $(x_n)$ is
uniformly named, its all-left root has value $u=\cut{\{x_n:n\in\N\}}{}$
and satisfies $u>x_n$ for every $n$. This proves the cofinality assertion
and the enumeration obstruction.
\end{proof}

The same argument gives a positive lower bound $\cut0{\{p_n:n\in\N\}}$
for every uniformly named family of positive values $p_n$. Neither operation
is analogous to taking an ordinary real supremum or infimum.

\subsection{A direct noncomputable real in the structural field}

\begin{proposition}[A halting real with a structural name]
\label{prop:haltingreal}
Let $K$ be the diagonal halting set. Then
\[
 r_K=\sum_{e\ge0}2\chi_K(e)4^{-(e+1)}
\]
is a noncomputable real with a valid structural name.
\end{proposition}
\begin{proof}
For each machine $e$, let a program initially enumerate no options. If
machine $e$ halts, it enumerates one left option with value zero. The resulting
name $h_e$ has value $0$ or $1$ according as $e\notin K$ or $e\in K$.
The names $h_e$ are uniform even though their values are not computable
uniformly as real numbers.

Using the arithmetic compilers, construct names for
$s_n=\sum_{e<n}2h_e4^{-(e+1)}$. Consider the structural cut
\begin{equation}\label{eq:haltingcut}
 \cut{\{s_n-4^{-n}:n\in\N\}}
      {\{s_n+4^{-n}:n\in\N\}}.
\end{equation}
Since $0\le r_K-s_n\le\tfrac23\,4^{-n}$, every left bound is strictly
below $r_K$ and every right bound strictly above it. The shrinking real
intervals determine every finite dyadic comparison with $r_K$. Their
simplest surreal separator is therefore the canonical real $r_K$: any
other separator has all the finite sign information of $r_K$ and is no
simpler. This is also the standard embedding of a real by its rational cuts.

Finally, the base-four digits of $r_K$ are only $0$ and $2$. Once the first
$e$ digits are known, the two possibilities for the next digit lie in
separated intervals; the tail after that digit is at most one third of its
nonzero contribution. Successively accurate Cauchy approximations would
therefore decide all of $K$. Thus $r_K\notin\Rc$.
\end{proof}

The general halting-real phenomenon is already present in Hamkins's
presentation \cite{hamkins}. The explicit margins in
\eqref{eq:haltingcut} are included to avoid the false assumption that a cut
may use an option equal to its proposed value.

\begin{corollary}[A structural--numerical incompatibility]
A notion which includes these uniform hereditary cut constructions and
Conway arithmetic cannot also have exactly $\Rc$ as its real slice. In
particular, there is no conversion of every real-valued structural name
into an ordinary computable Cauchy name.
\end{corollary}
\begin{proof}
Such a conversion would apply to \eqref{eq:haltingcut} and compute $r_K$.
\end{proof}


\begin{proposition}[Birthday and ordinal boundary]\label{prop:ordinal-boundary}
Every structurally computable surreal has birthday less than $\CK$.
The ordinals belonging to $\Sc$ are exactly the computable ordinals.
\end{proposition}
\begin{proof}
Unfold a valid enumerable dependency graph into a recursive tree of paths
carrying finite witnesses for each enumerated edge. Its well-founded
Kleene--Brouwer order is a computable well-order, so its rank is below $\CK$
\cite{sacks}. Induction bounds the birthday of each cut value by the rank
of its construction. The birthday of an ordinal is that ordinal itself.
Combine this bound with Proposition~\ref{prop:structcount}.
\end{proof}
This is a necessary complexity bound, not a characterization by birthday:
all ordinary real numbers have birthday at most $\omega$, but not all are
hyperarithmetic. Nor is a finite program carrying a semantic well-foundedness
promise the same as an algorithm which can decide that promise.

For uniformly given Boolean-valued valid structural names $b_n$, the names
$1-b_n$ and $b_nb_m$ implement negation and conjunction, while
$\cut{b_n-1:n\in\N}{}$ implements countable disjunction. Its value is zero
if every $b_n=0$, and one otherwise. The empty right side matters: putting
$1$ on that side would give $1/2$ in the true case. This elementary mechanism
helps explain the imported hyperarithmetic classification without purporting
to reprove it.

\subsection{Validity is a genuinely higher-order promise}

\begin{proposition}[Validity has full well-foundedness complexity]
\label{prop:validity}
For the program coding in Definition~\ref{def:structural}, the set of valid
root indices is $\Pi^1_1$-complete under computable many-one reductions.
\end{proposition}
\begin{proof}
For hardness, start with a recursive tree $T\subseteq\N^{<\N}$. At a node
$\sigma\in T$, put the programs for its immediate successors on the left
and nothing on the right. If $T$ is well founded, every cut is automatically
separated and the root is valid. If $T$ has an infinite branch, the dependency
graph is not well founded. The reduction is effective. Well-foundedness of
recursive trees is $\Pi^1_1$-complete \cite{sacks}.

For the upper bound, well-foundedness of the reachable dependency graph is
$\Pi^1_1$. On a well-founded graph, the recursive comparison relation on
pairs of positions is uniquely determined by
\begin{equation}
 a\preccurlyeq b\ \Longleftrightarrow\
 \bigl(\forall c\in L_a\;\neg(b\preccurlyeq c)\bigr)
 \ \land\
 \bigl(\forall d\in R_b\;\neg(d\preccurlyeq a)\bigr).
 \label{C-eq:comparison-rec}
\end{equation}
This recursion is well founded on pairs of ranks. Existence and uniqueness
hold even before knowing that every position is a number: it is the usual
comparison recursion for well-founded games. The assertion that a set of
pairs satisfies \eqref{C-eq:comparison-rec} is arithmetical in that set and the
program index. Hereditary separation is then the arithmetical requirement
that, at every reachable node, every left option is strictly below every
right option according to this relation.

Thus, in addition to well-foundedness, require that \emph{every} set of pairs
satisfying the recursive comparison equations satisfies hereditary
separation. This is a universal second-order quantification over an
arithmetical predicate. On a well-founded graph the unique comparison
relation exists, so the assertion is equivalent to validity and is not
vacuous. This gives a $\Pi^1_1$ definition.
\end{proof}

In particular there is no algorithm that recognizes exactly the valid
hereditary programs and halts on all of them. A proof-carrying subclass can
be effectively checked in a chosen proof system, but cannot exhaust the
semantic validity predicate.


\section{A computable order presentation can hide a noncomputable real}
\label{sec:signtrap}

A tempting alternative is to call a surreal computable when its signs are a
computable labeling of a computably presented well-order. The phrase sounds
like direct sign access, but it need not provide access to the sign at the
$n$th ordinal position. A computable presentation of an order of type
$\omega$ need not have a computable order isomorphism with the standard
$\N$.

\begin{theorem}[The signed-order presentation trap]
\label{thm:signtrap}
For every computably enumerable set $A\subseteq\N$, there is a decidable
domain $D_A\subseteq\N^2$, a decidable order on $D_A$ of order type
$\omega$, and a computable labeling by $+$ and $-$, such that the resulting
canonical sign sequence denotes a nondyadic real $r_A$ of Turing degree
exactly that of $A$.
\end{theorem}
\begin{proof}
Choose a computable stage enumeration in which an element enters $A$ at
most once. For every $e$, put $(e,0)$ and $(e,1)$ in $D_A$. Also put
$(e,s+2)$ in $D_A$ exactly when $e$ first enters $A$ at stage $s$.
Membership is decidable by simulating the enumeration through the specified
finite stage. Order $D_A$ lexicographically. Each $e$ contributes a nonempty
finite block of size two or three, so the order type is $\omega$.

Label $(e,0)$ by $+$ and every other point by $-$. The sign sequence in
ordinal order is the concatenation
\begin{equation}
 +\,-^{\,1+\chi_A(0)}\quad
 +\,-^{\,1+\chi_A(1)}\quad
 +\,-^{\,1+\chi_A(2)}\quad\cdots.
 \label{C-eq:block-signs}
\end{equation}
Both signs occur infinitely often. Its first two signs are $+-$, and its
successive dyadic intervals shrink to a nondyadic real in $(0,1)$.
With oracle $A$, one computes the blocks and hence obtains an ordinary fast
real name. Therefore $r_A\leq_T A$.

Conversely, a fast name for a nondyadic real computes any prescribed finite
initial segment of its surreal sign sequence. At each step compare with the
relevant dyadic pivot. Equality never occurs, and eventually the approximation
interval decides the strict comparison. After recovering sufficiently many
signs of \eqref{C-eq:block-signs}, count the one or two minus signs following
the $e$th plus sign. This recovers $\chi_A(e)$. Hence $A\leq_T r_A$.
\end{proof}

The theorem uses a particularly tame order: all blocks have size at most
three. What fails is not decidability of comparison or boundedness of block
size, but effective identification of the successor after a possibly missing
point. Thus one cannot establish real compatibility merely by placing
computability requirements on an abstract order and its labels.

For ordinary reals, \emph{canonical} sign access and Cauchy computation agree
at the level of computable values, with the familiar nonuniform treatment
of terminating dyadic expansions. For transfinite sign sequences, a proposed
representation must also specify ordinal addresses, restrictions to initial
segments, and the navigation information needed by its algorithms. The
coefficient representation developed next avoids this address problem
by supplying explicit rational coefficient addresses and effective finite search data.


\section{Why computable Hahn coefficients are insufficient}
\label{sec:badproduct}

Suppose we ask only that $q\mapsto a_q$ be computable, and that its nonzero
support be well ordered. Mathematical Hahn multiplication is defined
because each coefficient receives only finitely many contributions. It
does not follow that one can find those contributions or know that the
search for them is complete.

The next construction strengthens this objection. The supports are not
merely well ordered: they are decidable, left-finite, and have computably
bounded numbers of points below every exponent threshold.

\begin{theorem}[Decidable supports with a noncomputable product]
\label{thm:badproduct}
There are two rational-exponent Hahn series $A,B$, each with coefficients
in $\{0,1\}$ and a computable exact coefficient function on $\Q$, whose
supports are decidable and left-finite, but whose product does not have a
uniformly Cauchy-computable coefficient function.
\end{theorem}
\begin{proof}
Let $T(e)$ be the first halting stage when $e\in K$. Put
\[
 c_e=4^{e+1},\qquad \delta_s=\frac1{s+2},
\]
and define
\begin{equation}\label{eq:badAB}
 A=\sum_{e\in K}t^{c_e+\delta_{T(e)}},\qquad
 B=\sum_{e\in K}t^{2c_e-\delta_{T(e)}}.
\end{equation}
Each support has at most one point in its $e$th block. The $A$ blocks lie
in $(c_e,c_e+\tfrac12]$ and the $B$ blocks in
$[2c_e-\tfrac12,2c_e)$. They are separated and tend to infinity, proving
left-finiteness and a computable bound on the number of points below any
rational threshold.

For exact coefficient access to $A$, given $q\in\Q$, determine whether it
lies in one of the finitely searchable intervals
$(c_e,c_e+\tfrac12]$. If it does, compute
$s=1/(q-c_e)-2$. A nonnegative integral value $s$ is necessary; it is then
sufficient to simulate machine $e$ for $s$ stages and ask whether its first
halt was exactly at that stage. All other cases give coefficient zero.
The same algorithm for $B$ uses $s=1/(2c_e-q)-2$. This proves computability
of the coefficient functions and decidability of the supports.

We claim
\begin{equation}\label{eq:productencodes}
 \coef{3c_e}{(AB)}=\chi_K(e).
\end{equation}
A contributing pair indexed by $i,j\in K$ must satisfy
\[
 c_i+2c_j-3c_e=\delta_{T(j)}-\delta_{T(i)}.
\]
The left side is an integer divisible by four; the right side has absolute
value less than $1/2$. Hence both sides vanish. Uniqueness of finite
base-four expansions in
\[
 4^{i+1}+2\cdot4^{j+1}=3\cdot4^{e+1}
\]
forces $i=j=e$: distinct indices would produce a digit $1$ and a digit $2$
in different positions, rather than a single digit $3$. When $e\in K$,
that one pair is present and its perturbations cancel. Otherwise no pair
exists. This proves \eqref{eq:productencodes}.

A joint Cauchy algorithm for the product coefficients, requested to error
less than $1/4$ at $3c_e$, would distinguish zero from one and decide $K$.
The product is nevertheless an ordinary well-defined left-finite Hahn
series. Its failure is effective, not algebraic.
\end{proof}

The obstruction is not the impossibility of comparing real coefficients:
all input coefficients are exact integers. Nor is it the absence of a
cardinality bound on the finite work. The missing information is a
\emph{complete finite list of candidate exponents}. A bound saying ``at
most one point in this block'' does not reveal whether the point exists or
which rational encodes its halting time.

\begin{corollary}
The class of left-finite Hahn series with uniformly computable coefficient
functions is not, by itself, a field or even a ring under Hahn multiplication.
Decidable support membership and computable bounds on finite-section
cardinalities do not repair this failure.
\end{corollary}

This motivates a stronger numerical representation, but one which still
avoids requiring an impossible exact zero test on real coefficients.

\subsection{A complementary obstruction with singleton fibers}
The preceding construction keeps \emph{both} supports left-finite. The next,
retained from C, instead has a global bound of one on every convolution
fiber; one factor has a bounded accumulation point. These are distinct
strengthenings of the same negative principle, so neither is discarded.
\begin{theorem}[Sparse Hahn multiplication can compute the halting set]
\label{thm:singlefiber}
There are Hahn series $F,G$ with computable coefficient functions taking values in $\{0,1\}$ and
decidable well-ordered supports, such that $FG$ has no computable coefficient function. Every convolution
fiber for these two supports has size at most one.
\end{theorem}
\begin{proof}
For a halting machine $e$, write $s_e\geq0$ for its exact halting stage.
Define
\begin{equation}
 F=\sum_{n\geq0}t^{\,1-1/(n+2)},\qquad
 G=\sum_{e\in K}t^{\,e+1/(s_e+2)}.
 \label{C-eq:hahn-FG}
\end{equation}
The support of $F$ is strictly increasing, of order type $\omega$.
Its coefficient at $q=1$ is zero; at any other rational $q$, test whether
$(1-q)^{-1}$ is an integer at least two. The support of $G$ contains at most
one element in each interval $(e,e+1/2]$, so is also well ordered.
For a rational query $q$, let $e=\lfloor q\rfloor$. Test whether $e\geq0$
and $q-e=1/(s+2)$ for an integer $s\geq0$. If so, simulate machine $e$
through stage $s$ and test whether this is its exact halting stage.
This decides the coefficient of $G$. Thus both coefficient functions and
both supports are decidable.

An exponent in the product has the form
\[
 j+1+\frac1{s_j+2}-\frac1{n+2}.
\]
For fixed target exponent $q$, the integer $j$ must lie in
$(q-3/2,q-1/2)$, since $1-1/(n+2)\in[1/2,1)$ and
$1/(s_j+2)\in(0,1/2]$. This open interval of length one contains at most one integer. For each $j$ there is at most one support
term of $G$, and then at most one possible $n$. Thus every fiber has at
most one member. The ordinary Hahn product is well defined.

At the integer target $q=e+1$, the only possible integer block is $j=e$:
the difference $1/(s_j+2)-1/(n+2)$ is strictly between $-1/2$ and $1/2$,
so it cannot compensate a nonzero integer. The target is attained exactly
when $e$ halts and $n=s_e$. Consequently
\begin{equation}
 [t^{e+1}](FG)=\chi_K(e).
 \label{C-eq:product-halting}
\end{equation}
A uniform approximation algorithm for the coefficients of $FG$, requested
at these targets with error at most $1/4$, would decide $K$.
Therefore $FG$ has no computable coefficient function.
\end{proof}

The statement is stronger than nonexistence of a uniform multiplication
algorithm: the product of these two fixed computable coefficient functions
is not a computable coefficient function at all. Each individual coefficient
in \eqref{C-eq:product-halting} is of course a computable rational. It is their
uniform recovery that fails.

Classical support lemmas establish that each convolution coefficient is a
finite sum. They do not produce an effective stopping criterion for locating
its summands. Even a known global bound of one on their number does not
produce such a criterion. This is the same general danger as replacing an
unknown finite search domain by a promise that only finitely many answers
exist.


B uses the equivalent collision template
$\alpha_s=1-2^{-s-1}$,
$F_B=\sum_s t^{\alpha_s}$, and
$G_B=\sum_{e\in K}t^{2e+2-\alpha_{s_e}}$.
At target $2e+2$, equality of the two fractional offsets forces the same
halting stage, and the coefficient is $\chi_K(e)$. Rational queries recover
and check only a finite candidate stage. The singleton-fiber formulation
above and the left-finite formulation preceding it retain the stronger
conclusions; the dyadic offsets are an alternative fixture, not a
conflicting definition of multiplication.

\subsection{The sharp oracle cost of an unrestricted convolution}
\begin{theorem}[One jump for coefficients, not for an ordinary output name]
\label{thm:convolution-jump}
Let two Hahn series with rational exponents have uniformly computable real
coefficient functions and well-ordered supports, but no effective support
covers. Their product coefficient function is uniformly computable in $0'$.
The bound is sharp, already for exact rational zero--one coefficients.
\end{theorem}
\begin{proof}
For a fixed rational target $q$, the set
\[
 E_q=\{r\in\Q:a_r b_{q-r}\ne0\}
\]
is computably enumerable, since a nonzero computable real is semidecidable.
It is finite: infinitely many such $r$ would yield a strictly increasing
sequence in the first well-ordered support and hence a strictly decreasing
sequence $q-r$ in the second. This contradicts well-ordering.

Keep a finite list $D$ of contributors already detected. A halting oracle
decides whether the semidecision search for another member of $E_q\setminus D$
will succeed. If it will, run that search until a witness appears and add it;
otherwise stop. Finiteness ensures termination with exactly $E_q$.
Compute the finite sum of the named real products to any requested precision.
For exact rational coefficients this produces an exact rational result.
Either counterexample above gives product coefficients which decide $K$,
so any oracle uniformly solving this task computes $0'$.
\end{proof}
The witness search need not find a \emph{least} rational contributor and does
not require a decidable support test. This formulation also covers
Cauchy-real coefficients correctly. The resulting coefficient algorithm may
need $0'$ anew at every target $q$; in general no ordinary computable product
name exists. This differs from the finite normalization advice in
Theorem~\ref{thm:jump}, after which an ordinary inverse program is returned.
The theorem supplies coefficients, not an effective finite-cover name for
the product. Structured positive results for special Hahn families, such as
Mahler-equation support receptacles \cite{faverjon}, impose additional data
and do not contradict this obstruction.

\section{Effective left-finite Cauchy--Hahn names}
\label{sec:definition}

\begin{definition}[The numerical representation]\label{def:elf}
An \emph{effective left-finite Cauchy--Hahn name} consists of a rational
number $\ell$ and two total computable procedures $A,F$ with the following
meaning. There is a real coefficient family $(a_q)_{q\in\Q}$ such that
\begin{enumerate}[label=\textup{(\arabic*)},leftmargin=2em]
\item $a_q=0$ whenever $q<\ell$;
\item $A(q,k)\in\Q$ and $|A(q,k)-a_q|\le2^{-k}$ for every $q,k$;
\item for each $B\in\Q$, $F(B)$ is a finite rational set contained in
$[\ell,B)$ and containing every $q<B$ with $a_q\ne0$.
\end{enumerate}
Its value is $x=\sum_qa_qt^q$, viewed in $\No$ through normal forms.
The set $\Lc$ consists of the values of such computable names.
\end{definition}

The procedure $F$ is a \emph{finite-section cover}, not an exact support
oracle. It may list exponents with coefficient zero, and no test for deleting
such entries is required. Lists can always be converted effectively into
finite sets of reduced rational numbers. They need not be nested as $B$
increases. For $B\le\ell$ the only allowed list is empty.

These small details matter. Requiring an exact nonzero-support list would
already demand a zero test for a single arbitrary computable real constant.
Allowing harmless zero candidates makes addition and cancellation effective.
The name's correctness conditions remain promises; the algorithms below
preserve them rather than attempting to decide them.

\begin{proposition}[Well-definedness and the real slice]
\label{prop:realslice}
Every valid numerical name defines a unique left-finite Hahn series.
The set $\Lc$ is countable, contains every rational monomial $t^q$, and
satisfies $\Lc\cap\R=\Rc$. The inclusion of ordinary computable real
constants into this representation is uniform.
\end{proposition}
\begin{proof}
The error bounds determine each real $a_q$ uniquely. The covers imply
left-finiteness; hence the series is a valid Hahn normal form. Its surreal
value is unique by the normal-form theorem. Countably many programs and
finite rational data give countably many names.

For $t^q$, use lower bound $q$, coefficient one at $q$ and zero elsewhere,
and cover $\{q\}$ exactly when $B>q$. For a Cauchy-named real $a$, use
lower bound zero, coefficient $a$ at exponent zero and zero elsewhere,
and cover $\{0\}$ when $B>0$. This construction is uniform even when
$a=0$. Conversely, if the value of a name is an ordinary real, uniqueness
of normal form says that its only possible nonzero coefficient is $a_0$.
The algorithm $A(0,k)$ therefore computes that real.
\end{proof}

\begin{definition}[Finite jet]
For $B\in\Q$ define the exact finite jet
\[
 J_Bx=\sum_{q\in F(B)}a_qt^q.
\]
Its data consist of a finite rational exponent list and Cauchy names for its
real coefficients, not finite rational approximations to those coefficients.
\end{definition}

\begin{proposition}[Certified truncation]
\label{prop:jets}
From a numerical name and $B$ one computes a name for $J_Bx$, and
\[
 \val(x-J_Bx)\ge B.
\]
The value of $J_Bx$ is independent of the particular valid cover.
\end{proposition}
\begin{proof}
All nonzero coefficients below $B$ occur in the list, and the remaining
listed coefficients are zero. Thus the coefficients of the difference
below $B$ vanish. Restrict the original coefficient oracle to the finite
list to obtain a name for the jet.
\end{proof}

The series in Theorem~\ref{thm:badproduct} have no computable covers of this
kind. For example, a cover of $A$ below $c_e+1$ would give finitely many
candidates in its $e$th block. Its exact coefficient oracle could inspect
them and decide whether $e\in K$.

\section{Uniform ring operations and certified reciprocation}
\label{sec:arithmetic}

\subsection{Finite Cauchy arithmetic}

We repeatedly use the elementary fact that finite sums and products of
Cauchy-computable real numbers are computable uniformly. This does not
require zero tests. A precision-zero approximation gives a computable bound
$|a|\le |A(0)|+1$. To approximate a finite sum of products, choose a common
input error $\eta\le1$ small enough that the sum of the estimates
\[
 |ab-\widetilde a\widetilde b|
 \le \eta\bigl(|a|+|b|+1\bigr)
\]
is below the requested output error. All required bounds follow from coarse
approximations. If $a\ne0$ is promised, its reciprocal is computable as an
ordinary real: search until an approximation interval excludes zero and
then propagate errors through division using the resulting positive lower
bound for $|a|$.

\begin{theorem}[Uniform numerical ring arithmetic]
\label{thm:ring}
Addition, negation, and multiplication are computable uniformly on
numerical names. For rational $r$, the operation $x\mapsto t^rx$ is also
uniformly computable.
\end{theorem}
\begin{proof}
Negation keeps the lower bound and cover. For addition use lower bound
$\min(\ell_x,\ell_y)$, coefficient sums, and the cover
$F_x(B)\cup F_y(B)$. Rational shifting adds $r$ to the lower bound and
exponents, uses $a_{q-r}$ as coefficient, and shifts $F_x(B-r)$ by $r$.

For multiplication put $\ell=\ell_x+\ell_y$. At an exponent $q$ the exact
coefficient is the finite sum
\begin{equation}\label{eq:coeffmul}
 \coef q{(xy)}=
 \sum_{r\in F_x(q-\ell_y+1)} (\coef rx)(\coef{q-r}y).
\end{equation}
Indeed, a nonzero contribution requires $r\ge\ell_x$ and
$q-r\ge\ell_y$, so $r\le q-\ell_y$ and lies in the indicated cover.
Zero candidates contribute nothing. Finite Cauchy arithmetic computes this
coefficient to any requested precision.

A support cover for the product below $B$ is
\begin{equation}\label{eq:covermul}
 F_{xy}(B)=
 \bigl(F_x(B-\ell_y)+F_y(B-\ell_x)\bigr)\cap[\ell,B),
\end{equation}
where the plus sign between finite sets denotes their pairwise sum set.
If $r+s<B$ contributes to a nonzero coefficient, then
$r<B-\ell_y$ and $s<B-\ell_x$. Thus \eqref{eq:covermul} contains every
possible nonzero exponent below $B$. It is a computable finite set. This
proves all requirements on the product name.
\end{proof}

No test for cancellation was hidden in either formula. In particular, the
output may have a much larger true leading exponent than its supplied
lower bound, or may be zero.

\subsection{An effective positive gap}

\begin{lemma}[Gap extraction without normalization]
\label{lem:gap}
Suppose a numerical name is promised to satisfy
$\supp(h)\subseteq\Q_{>0}$. One can compute $\delta\in\Q_{>0}$ such that
$\val(h)\ge\delta$, including the case $h=0$.
\end{lemma}
\begin{proof}
Compute $F_h(1)$, discard its nonpositive candidates, and take
\[
 \delta=\min\bigl(\{1\}\cup(F_h(1)\cap\Q_{>0})\bigr).
\]
Every nonzero exponent of $h$ below one occurs among these positive
candidates; every other nonzero exponent is at least one. The promise
justifies discarding nonpositive candidates, even though their zero
coefficients were not algorithmically recognized.
\end{proof}

Under the same promise we may replace the lower bound by zero and discard
all nonpositive candidates in every cover. This is \emph{certified
filtering}: a theorem about the input, rather than an equality decision,
authorizes removal of those positions.

\begin{lemma}[Effective geometric series]
\label{lem:geom}
For a numerically named $h$ with positive support, the series
$G(h)=\sum_{n\ge0}h^n$ has a uniformly computable numerical name. It satisfies
$(1-h)G(h)=1$ and, for every $N\ge0$,
\[
 (1-h)\sum_{n=0}^{N}h^n=1-h^{N+1}.
\]
\end{lemma}
\begin{proof}
Extract a positive gap $\delta$. Each $h^n$ is numerically computable by
Theorem~\ref{thm:ring} and has support in $[n\delta,\infty)$. At any
threshold $B$, only the finitely many $n$ with $n\delta<B$ can contribute.
The union of their computed covers is a finite cover for $G(h)$ below $B$.
For a queried exponent $q$, choose $N$ with $(N+1)\delta>q$ and compute
the coefficient of the finite sum through $N$. This stabilizes the exact
coefficient, after which its real precision is computed separately.
A lower bound is zero.

Finite multiplication proves the displayed remainder identity. At every
fixed coefficient, $h^{N+1}$ eventually contributes zero. Hence the
coefficientwise Hahn identity $(1-h)G(h)=1$ follows. This reasoning takes
place in the fixed left-finite workspace, not in the full surreal topology.
\end{proof}

\begin{theorem}[Field closure and certified inversion]
\label{thm:inverse}
Every nonzero $x\in\Lc$ has $x^{-1}\in\Lc$. Uniform inversion is possible
from a numerical name supplemented by its true leading exponent
$\lambda=\val(x)\in\Q$.
\end{theorem}
\begin{proof}
The coefficient $a=\coef\lambda x$ is a nonzero computable real. With the
nonzero promise compute its real inverse. Form
\[
 h=1-a^{-1}t^{-\lambda}x.
\]
The supplied leading-exponent certificate proves that all coefficients of
$h$ at nonpositive exponents vanish. Lemma~\ref{lem:geom} now gives
\begin{equation}\label{eq:inverse}
 x^{-1}=a^{-1}t^{-\lambda}\sum_{n\ge0}h^n
\end{equation}
uniformly, with lower bound $-\lambda$ and covers obtained by shifting
those of the geometric series. The exact identity follows from
$1-h=a^{-1}t^{-\lambda}x$.

For extensional closure, every particular nonzero left-finite series has a
least support exponent, which is a rational number. That single rational
can be hard-coded into an inverse program. This proves existence of a
computable inverse name without claiming a uniform procedure for finding
that rational from an arbitrary input name.
\end{proof}

This final distinction is the source of the apparent paradox: all inverse
values are computable in the specified sense, but one cannot uniformly
choose their names from all nonzero input names. The obstruction will be
proved, rather than merely asserted, in Section~\ref{sec:uniformlimits}.

\subsection{The valuation ring and standard part}

Define
\[
 \Oh=\{x\in\Lc:\val(x)\ge0\},\qquad
 \mh=\{x\in\Lc:\val(x)>0\}.
\]
These include zero by the convention $\val(0)=+\infty$.

\begin{proposition}[Value group and effective residue]
\label{prop:residue}
The value group of $\Lc$ is $\Q$. The ring $\Oh$ is convex in its ordered
field and has maximal ideal $\mh$. Coefficient extraction gives
\[
 \st:\Oh\longrightarrow\Rc,\qquad \st(x)=\coef0x,
 \qquad \Oh/\mh\cong\Rc.
\]
This standard-part operation is computable from a numerical name with the
promise $x\in\Oh$, without knowing $\val(x)$.
\end{proposition}
\begin{proof}
Every nonzero valuation is a rational exponent, and $\val(t^q)=q$ shows
surjectivity onto $\Q$. An element has nonnegative valuation exactly when
it is finite as a surreal: a negative leading exponent yields magnitude
larger than every positive integer, and a nonnegative exponent gives a real
constant plus an infinitesimal. Thus $\Oh$ is a convex subring.

For two elements of $\Oh$, only their exponent-zero terms contribute to
the exponent-zero coefficient of their product. Hence $\st$ is a ring
homomorphism. Its kernel is $\mh$, and constants give surjectivity. The
quotient is a field, proving maximality. Finally, $A(0,k)$ is already an
ordinary Cauchy name for the residue.
\end{proof}

Each finite element has the unique decomposition
$x=\st(x)+(x-\st(x))$ with infinitesimal second summand. The algorithm
extracts it without deciding whether that infinitesimal summand is zero.

\section{Effective Hensel lifting and real closedness}
\label{sec:realclosed}

\subsection{A quantitative simple-root algorithm}

For a polynomial $P\in\Oh[Z]$, denote by $\overline P\in\Rc[Z]$ its
coefficientwise residue. A residue root $c$ is simple if
$\overline P(c)=0$ and $\overline{P'}(c)\ne0$.

\begin{theorem}[Effective simple Hensel lifting]
\label{thm:hensel}
Let a finite-degree polynomial $P\in\Oh[Z]$ be given by numerical names for
its coefficients, and let $c\in\Rc$ be given by a Cauchy name. Under the
promises
\[
 \overline P(c)=0,\qquad a=\overline{P'}(c)\ne0,
\]
one can compute a numerical name for the unique $y\in\Oh$ satisfying
$P(y)=0$ and $y-c\in\mh$.
More precisely, one computes a positive rational $\delta$ and numerical
Newton iterates $y_n$ such that
\begin{equation}\label{eq:newtonbounds}
 \val(P(y_n))\ge2^n\delta,\qquad
 \val(y-y_n)\ge2^n\delta.
\end{equation}
\end{theorem}
\begin{proof}
Put $y_0=c$. Since $P(c)$ has positive support, the gap lemma gives
$\delta>0$ with $\val(P(c))\ge\delta$. This also works if $P(c)=0$.
Inductively define
\begin{equation}\label{eq:newton}
 y_{n+1}=y_n-P(y_n)P'(y_n)^{-1}.
\end{equation}
If $y_n\equiv c\pmod\mh$, then $P'(y_n)$ has residue $a\ne0$ and hence
\emph{known} leading exponent zero. Its inverse is therefore computed by
the certified algorithm, without a general normalization procedure.
The correction has positive valuation, so $y_{n+1}\equiv c\pmod\mh$.

All polynomial coefficients and iterates lie in $\Oh$. The finite Taylor
identity, with $u=-P(y_n)/P'(y_n)$, gives
\[
 P(y_n+u)=P(y_n)+P'(y_n)u+u^2R_n,
 \qquad R_n\in\Oh.
\]
The first two terms cancel. If $\val(P(y_n))\ge2^n\delta$, the unit
$P'(y_n)$ has valuation zero, so $\val(u)\ge2^n\delta$ and
$\val(P(y_{n+1}))\ge2^{n+1}\delta$. Induction proves these estimates
and $\val(y_{n+1}-y_n)\ge2^n\delta$.

For each $q\in\Q$, choose $n$ with $2^n\delta>\max(q,0)$ and define the
coefficient of $y$ at $q$ to be that of $y_n$. All later iterates agree at
that exponent, so this is consistent. It gives a uniform Cauchy coefficient
algorithm. Set every coefficient at a negative exponent to zero, as
justified by the invariant $y_n\in\Oh$. For a threshold $B>0$, choose
$n$ with $2^n\delta>B$ and use $F_{y_n}(B)\cap[0,B)$ as a cover for $y$;
for $B\le0$, use the empty set. Thus $y$ has a numerical name with lower
bound zero, and the second inequality in \eqref{eq:newtonbounds} follows.

For any fixed bound, polynomial evaluation uses only finitely many earlier
coefficients, since all arguments and coefficients have nonnegative
valuation. Consequently $P(y)$ agrees below that bound with $P(y_n)$ for
sufficiently large $n$, and the first inequality in
\eqref{eq:newtonbounds} makes those coefficients zero. Thus $P(y)=0$.

For uniqueness, if $z\equiv c\pmod\mh$ is another root, the polynomial
difference quotient gives
\[
 P(y)-P(z)=(y-z)U,\qquad U\equiv P'(c)\equiv a\pmod\mh.
\]
The factor $U$ is a unit, so $y=z$.
\end{proof}

The number of Newton rounds needed to determine an exponent jet grows
logarithmically in $B/\delta$. This is a bound on rounds, not a polynomial
bit-complexity claim: cover sizes, coefficient magnitudes, and real
precision costs may grow substantially. The theorem is effective only
under its promises; it does not claim to decide whether a proposed residue
root is simple.

\subsection{Why the computable reals are real closed}

For completeness, the extensional real-algebraic fact needed below has an
elementary proof which does not require a uniform multiplicity decision.

\begin{lemma}\label{lem:Rcclosed}
Every real root of a nonzero polynomial with coefficients in $\Rc$ belongs
to $\Rc$. Consequently $\Rc$ is real closed.
\end{lemma}
\begin{proof}
Let $r$ be a real root of multiplicity $m$ of $P$. The polynomial
$Q=P^{(m-1)}$ has a simple root at $r$. Choose rational endpoints $u<r<v$
and a positive rational $d$ such that $Q'$ has constant sign and
$|Q'(x)|>d$ throughout $[u,v]$, and $r$ is the unique root there.
These finitely many data exist by continuity and simplicity and may be
hard-coded for this particular root.

For a requested precision $k$, search through rational $q\in(u,v)$ and
increasing computation precisions until
$|Q(q)|<d\,2^{-k}$ is certified. Polynomial evaluation at a rational is
Cauchy computable, and this strict inequality is semidecidable. Rationals
arbitrarily close to $r$ ensure that the search succeeds. The mean value
theorem then gives $|q-r|<2^{-k}$. Hence $r$ is computable.

Every real element algebraic over $\Rc$ is such a root. Since $\R$ is
real closed, positive square roots and roots of odd-degree polynomials
exist there and therefore already lie in $\Rc$. These two properties
characterize real-closed ordered fields.
\end{proof}

The proof chooses $m,u,v,d$ nonuniformly. This is sufficient for the
algebraic statement and should not be mistaken for a root-isolation
algorithm that determines multiplicities from arbitrary coefficient names.

\subsection{The one external valuation theorem}

We use the following standard consequence of the lemma of Ostrowski; the
formula and the residue-characteristic-zero specialization are recorded
explicitly in \cite[\S2.1, equation (1)]{rzepka}:
\begin{equation}\label{eq:ostrowski}
 [E:F]=(\val E:\val F)[Ev:Fv]
\end{equation}
for a finite extension of a henselian valued field $F$ of residue
characteristic zero, with its unique extended valuation. Here $Fv$ and
$Ev$ denote residue fields. The general formula has an additional defect
factor, a power of the residue characteristic exponent; in characteristic
zero that exponent is one. Henselianity, equivalently uniqueness of the
valuation extension, is essential.

\begin{lemma}[A real-closedness criterion]
\label{lem:valuationRC}
Let $F$ be an ordered field with a convex henselian valuation. If its value
group is divisible and its residue field is real closed of characteristic
zero, then $F$ is real closed.
\end{lemma}
\begin{proof}
Let $E/F$ be a finite ordered algebraic extension. The convex hull in $E$
of the valuation ring of $F$ is a convex valuation ring extending it.
Indeed, a convex subring of an ordered field is a valuation ring: if
$x$ does not lie in it, $|x|$ exceeds every positive element of the ring,
so $|x^{-1}|<1$ and $x^{-1}$ lies in it. Its intersection with $F$ is the
original convex ring. The induced residue field is ordered, extending the
order on $Fv$. Henselianity makes this the unique extended valuation.

The fundamental inequality makes both the value-group index and residue
degree finite. A torsion-free ordered group cannot properly contain a
divisible subgroup with finite index: for $\gamma\in\val E$, some
$n\gamma$ lies in $\val F$; divisibility supplies $\eta\in\val F$ with
$n\eta=n\gamma$, and torsion-freeness gives $\eta=\gamma$. Thus the
ramification index is one. The finite ordered algebraic extension $Ev/Fv$
is trivial because $Fv$ is real closed, so the residue degree is one.
Equation~\eqref{eq:ostrowski} gives $[E:F]=1$.

If $F$ were not real closed it would have a proper ordered algebraic
extension, containing a proper finite subextension generated by one element.
The preceding argument rules this out.
\end{proof}

\begin{theorem}[Real closedness of the effective core]
\label{thm:Lcrealclosed}
The ordered field $\Lc$ is real closed, and $\Lc[i]$ is algebraically
closed.
\end{theorem}
\begin{proof}
Field closure was proved in Theorem~\ref{thm:inverse}. The valuation is
convex, has value group $\Q$, and has residue field $\Rc$, by
Proposition~\ref{prop:residue}. Theorem~\ref{thm:hensel} gives the simple
root form of Hensel's lemma, hence henselianity. The group $\Q$ is divisible,
and Lemma~\ref{lem:Rcclosed} makes the characteristic-zero residue field
real closed. Apply Lemma~\ref{lem:valuationRC}. The usual algebraic
characterization of a real-closed field says that adjoining a square root
of $-1$ gives its algebraic closure.
\end{proof}

This proves more than the existence of square roots of infinitesimal
perturbations of one: \emph{every} odd-degree polynomial over the effective
core has a root in that core, regardless of the scales or multiplicities
of its roots. The general existence assertion is extensional. The uniform root algorithms proved here require stated residue data or
a finite-jet branch certificate; neither is a blind branch selector.

\begin{corollary}[Algebraic operations on individual computable surreals]
Every real surreal root of a \emph{nonzero} polynomial over $\Lc$ belongs to $\Lc$.
Every surcomplex root of such a polynomial belongs to $\Lc[i]$.
\end{corollary}
\begin{proof}
A real-closed ordered subfield has no proper ordered algebraic extension.
Any real surreal root generates such an extension unless it already lies
in the field. For a surcomplex root, the polynomial already splits over
$\Lc[i]$, which is algebraically closed, so every root in a larger field
is one of its roots there.
\end{proof}

\section{The computable Puiseux subfield}
\label{sec:puiseux}

A particularly transparent subfield uses a fixed denominator for all
exponents of each element, although different elements may have different
denominators.

\begin{definition}[Computable Puiseux series]
\label{def:puiseux}
Let $\Pc$ consist of all series
\begin{equation}\label{eq:Puiseux}
 x=\sum_{n\ge N}a_nt^{n/m},
 \qquad m\in\N_{>0},\quad N\in\Z,
\end{equation}
with uniformly Cauchy-computable coefficients $a_n\in\Rc$. A name consists
of $m,N$ and this joint coefficient algorithm.
\end{definition}

\begin{proposition}[Embedding and arithmetic]
\label{prop:puiseuxarithmetic}
Puiseux names translate uniformly to numerical names. The set $\Pc$ is a
subfield of $\Lc$. Addition and multiplication are uniform; inversion is
uniform when the leading integer index is supplied.
\end{proposition}
\begin{proof}
Use lower bound $N/m$. At a rational exponent $q$, return $a_{mq}$ if
$mq$ is an integer at least $N$, and zero otherwise. A cover below $B$ is
the finite grid $\{n/m:N\le n<mB\}$. For two series pass to denominator
$\lcm(m_1,m_2)$, inserting zero coefficients at the other grid points.
Finite convolution computes sums and products.

For a nonzero series with leading index $\nu$, write
$x=t^{\nu/m}\sum_{j\ge0}a_{\nu+j}t^{j/m}$. Its inverse coefficients are
computed successively from
\begin{equation}\label{eq:inverseRecurrence}
 b_0=a_\nu^{-1},\qquad
 b_n=-a_\nu^{-1}\sum_{j=1}^n a_{\nu+j}b_{n-j}\quad(n\ge1),
\end{equation}
so $x^{-1}=t^{-\nu/m}\sum_{n\ge0}b_nt^{n/m}$. Each requested coefficient
requires only a finite recursive computation with real Cauchy arithmetic.
For an individual nonzero value, its integer $\nu$ can be hard-coded.
\end{proof}

\begin{theorem}[A proper real-closed subfield]
\label{thm:puiseuxclosed}
The field $\Pc$ is real closed and is a proper subfield of $\Lc$.
\end{theorem}
\begin{proof}
For a finite polynomial over the valuation ring of $\Pc$, choose a common
exponent denominator $m$ for its coefficients. In the simple Hensel
construction, all Newton iterates and their inverses stay on this same
$\frac1m\Z$ grid. The stabilized coefficients are uniformly computable,
with lower index zero for the lifted integral root. Thus the proof of
Theorem~\ref{thm:hensel} gives henselianity inside $\Pc$, not merely inside
$\Lc$. Its value group is still $\Q$, since it contains every $t^q$, and
its residue field is $\Rc$. Lemma~\ref{lem:valuationRC} proves real
closedness.

For strictness consider
\begin{equation}\label{eq:unboundeddenoms}
 u=\sum_{n\ge1}t^{\,n+1/(n+1)}.
\end{equation}
The exponent sequence is computable, strictly increasing, and unbounded.
A coefficient query $q$ only requires testing the finitely many integers
$1\le n<q$ for the exact equation $q=n+1/(n+1)$. A cover below $B$ is
computed by checking $1\le n<B$. This gives a numerical name for $u$.
But the reduced denominator of $n+1/(n+1)$ is exactly $n+1$, since
$\gcd(n(n+1)+1,n+1)=1$. These denominators cannot all divide any fixed
integer $m$. Uniqueness of Hahn normal form shows $u\notin\Pc$.
\end{proof}

No Newton--Puiseux algorithm was silently assumed here. Real closedness
follows from the effective unit-root iteration plus an explicit standard
valuation criterion. General roots may require ramification or nonuniform
branch data; the theorem does not compute all those data blindly.

\begin{example}[Explicit computable infinitesimals]
All of
\[
 \frac{1}{1-t}=\sum_{n\ge0}t^n,\qquad
 e^t=\sum_{n\ge0}\frac{t^n}{n!},\qquad
 \sum_{n\ge0}\pi^{n^2}t^{n/3}
\]
belong to $\Pc$. There is no real-analytic growth condition on the
coefficients of a formal infinitesimal series: the last example is allowed
regardless of the radius of convergence of a related ordinary real power
series. Its grid supplies finite exponent sections.
\end{example}

The numerical core is deliberately not closed under all small Conway
cuts. Every element of $\Lc$ is eventually dominated in magnitude by some
$\omega^n$, because its leading exponent is rational. Thus the uniformly
named sequence $1,\omega,\omega^2,\ldots$ is cofinal in its positive
magnitudes. Its all-left Conway cut is $\omega^\omega$, which cannot lie
in $\Lc$: its normal form has exponent $-\omega$ in the $t$ convention,
not a rational exponent. Structural computability accommodates this cut;
our fixed rank-one numerical workspace does not. Larger scales require
larger workspaces, not a change in the meaning of rational exponents.

\section{Finite jets and real-closedness}
\label{C-sec:realclosed}

The arguments in B and C are the same finite-advice mechanism, written in
slightly different normalizations. We retain C's explicit divided-derivative
formulation and B's emphasis on a single fixed coefficient inverse.
It is not enough to invoke Newton--Puiseux and observe that every root
coefficient is a computable real. The coefficients of one root must have a
\emph{single} algorithm. This section proves that uniformity. The proof is
local at a simple root and needs only a finite amount of nonuniform seed
data.

For any coefficient field $k$, write
$\PP(k)=\bigcup_{d\ge1}k((t^{1/d}))$ without an effectiveness restriction.
Importantly, $\Pc\subsetneq\PP(\Rc)$: the series
$\sum_{n\ge0}\chi_K(n)t^n$ has individually computable rational coefficients
but no jointly computable coefficient function. In every Puiseux name,
coefficients at integer exponents are uniformly readable.

\subsection{The classical algebraic input}

We use the characteristic-zero Newton--Puiseux theorem in the following
standard form: if $k$ is algebraically closed of characteristic zero, then
$\PP(k)=\bigcup_d k((t^{1/d}))$ is algebraically closed. The introductions of Rond and Kedlaya give this formulation and its relation
to finite ramification \cite{rond,kedlaya}. This is the algebraic-existence input in this second route; the first
route in Section~\ref{sec:puiseux} uses henselian valuation theory instead.

It follows that $\PP(k)$ is real closed when $k$ is real closed. Indeed,
$k(i)$ is algebraically closed, and
\[
 \PP(k)(i)=\PP(k(i))
\]
by separating real and imaginary coefficients at a common denominator.
The field $\PP(k)$ is ordered by its leading coefficient, hence formally
real, and adjoining $i$ makes it algebraically closed. Apply this with
$k=\Rc$, using Lemma~\ref{lem:Rcclosed}.

This classical statement concerns arbitrary coefficient sequences. The
following lemma is the additional effective step.

\subsection{A coefficientwise effective simple-root lemma}

For a polynomial in $Y$, write $f^{[j]}=f^{(j)}/j!$ for its divided formal
derivative. Thus
\[
 f(Y+H)=\sum_{j\geq0}f^{[j]}(Y)H^j,
\]
where the sum has only finitely many terms.

\begin{lemma}[Finite-jet lifting]
\label{C-lem:finite-jet}
Let $f(Y)\in\Rc[[u]][Y]$ have uniformly computable coefficients as power
series in $u$. Suppose $y\in\Rc[[u]]$ is a root, with
$f'(y)\ne0$. Then the coefficients of $y$ are uniformly computable.
More explicitly, they can be computed uniformly from $f$ and the finite
seed data
\[
 s=v_u(f'(y)),\qquad M>s,\qquad
 a=\sum_{j=0}^{M-1}y_j u^j,
\]
where the finitely many coefficients of $a$ are supplied by fast real names.
\end{lemma}
\begin{proof}
Since $f$ and $y$ are integral power series, $s$ is a nonnegative integer.
Write $y=a+u^Mz$ with $z\in\Rc[[u]]$.
For integral arguments, polynomial evaluation is valuation nondecreasing
with respect to their difference. In particular,
\[
 f'(a)-f'(y)\in u^M\Rc[[u]].
\]
As $M>s$, this implies $v_u(f'(a))=s$.
Taylor expansion about the actual root gives
\[
 f(a)=f'(y)(a-y)+\sum_{j\geq2}f^{[j]}(y)(a-y)^j.
\]
The first term has valuation at least $M+s$; every other term has valuation
at least $2M>M+s$. Therefore $v_u(f(a))\geq M+s$.

Consider the transformed polynomial
\begin{equation}
 Q(Z)=u^{-(M+s)}f(a+u^MZ).
 \label{C-eq:Q-transform}
\end{equation}
Its constant coefficient is integral by the preceding bound. Its linear
coefficient is $u^{-s}f'(a)$, an integral unit. For $j\geq2$, the coefficient
of $Z^j$ is
\[
 u^{(j-1)M-s}f^{[j]}(a),
\]
which is divisible by $u$ because $M>s$. Thus $Q\in\Rc[[u]][Z]$, and its
reduction modulo $u$ is
\begin{equation}
 Q(Z)\bmod u=b+cZ,\qquad c\in\Rc\setminus\{0\}.
 \label{C-eq:Q-mod}
\end{equation}
All coefficients of $Q$ are uniformly computable from $f$ and the finite
seed data: the transformation consists of finite polynomial operations,
coefficient shifts, and finite real arithmetic.

There is a unique formal root $z=\sum_{n\geq0}z_nu^n$ of $Q$ in
$\Rc[[u]]$. Its constant term must be $z_0=-b/c$. Once
$z_0,\ldots,z_{n-1}$ are known, the coefficient of $u^n$ in $Q(z)$ has the
form
\begin{equation}
 c z_n+E_n(z_0,\ldots,z_{n-1}),
 \label{C-eq:triangular}
\end{equation}
where $E_n$ is a computably obtainable finite polynomial expression in those
coefficients and finitely many coefficients of $Q$. No nonlinear term can
introduce an additional occurrence of $z_n$ at degree $n$, since all the
nonlinear coefficients of $Q$ are divisible by $u$. Higher terms of its
linear coefficient also use only earlier $z_j$.

Set $z_n=-c^{-1}E_n$ successively. Reciprocation of the nonzero computable
real $c$ and every finite arithmetic operation are computable. Recursion on
$n$ therefore supplies one algorithm for $(z_n)$, and uniqueness follows at
the same time from \eqref{C-eq:triangular}. The previously given
$(y-a)/u^M$ is a root, so it is this computed series. Appending the finite
prefix $a$ recovers a uniform coefficient algorithm for $y$.
\end{proof}

\begin{remark}[Why the finite seed is legitimate]
The hypothesis $y\in\Rc[[u]]$ says only that each coefficient individually
is a computable real. It does not supply a uniform algorithm for the whole
sequence. However, finitely many such coefficients can be equipped with
finitely many programs, all hard-coded into one algorithm. Lemma~\ref{C-lem:finite-jet}
uses exactly this finite information to compute every remaining coefficient.
It does not hard-code infinitely many programs and does not assume the
conclusion it is proving.
\end{remark}

\begin{corollary}[Laurent version]
\label{C-cor:laurent-lift}
If $f\in\Pc[Y]$ and $y\in\PP(\Rc)$ is a simple root of $f$, then
$y\in\Pc$.
\end{corollary}
\begin{proof}
Choose a common finite ramification denominator for the finitely many
coefficients of $f$ and for $y$, and write its coordinate as $u$.
Choose an integer $h\geq0$ making $z=u^hy$ integral. Replacing $Y$ by
$u^{-h}Z$ in $f$ and multiplying the polynomial by a sufficiently large
power of $u$ gives $g(Z)\in\Rc[[u]][Z]$ with uniformly computable
coefficients. The root $z$ remains simple. Apply
Lemma~\ref{C-lem:finite-jet}, then shift back by $u^{-h}$.
All choices made here are finite data for the fixed polynomial and root.
\end{proof}

\subsection{Algebraic closure inside the full Puiseux field}

\begin{theorem}[Real-closedness of the operational model]
\label{C-thm:pc-realclosed}
The field $\Pc$ is relatively algebraically closed in $\PP(\Rc)$.
Consequently $\Pc$ is a real-closed ordered subfield of $\No$.
Moreover, $\Pc[i]$ is algebraically closed.
\end{theorem}
\begin{proof}
Let $y\in\PP(\Rc)$ be algebraic over $\Pc$. Choose its minimal polynomial
$f\in\Pc[Y]$. This makes sense because Proposition~\ref{prop:puiseuxarithmetic} has
already established that $\Pc$ is a field. Characteristic zero makes the
minimal polynomial separable, so $y$ is a simple root. The finitely many
coefficients of $f$ admit a uniformly combined finite family of Puiseux
coefficient names. Corollary~\ref{C-cor:laurent-lift} yields $y\in\Pc$.
This proves relative algebraic closedness.

The ambient field $\PP(\Rc)$ is real closed by the classical input above.
A positive $x\in\Pc$ has its positive square root in that ambient field;
an odd-degree polynomial over $\Pc$ has a real root there. Each such root
is algebraic over $\Pc$ and hence belongs to $\Pc$ by the first part.
This is the real-closedness criterion. Finally, the standard
characterization of real-closed fields says that adjoining a square root
of $-1$ yields an algebraically closed field.
\end{proof}

\begin{corollary}
Every real root in $\No$ of a nonzero polynomial with coefficients in $\Pc$
is in $\Pc$. Every root in $\No[i]$ of a nonzero polynomial with coefficients
in $\Pc[i]$ is in $\Pc[i]$.
\end{corollary}
\begin{proof}
An ordered extension of a real-closed field has no proper algebraic ordered
subextension, and an algebraically closed field has no proper algebraic
extension. Alternatively, factor the polynomial over $\Pc[i]$: its full
finite list of linear factors leaves no additional roots in a larger field.
\end{proof}

The theorem supplies much more than closure under radicals. For example,
all five roots in $\No[i]$ of
\[
 Y^5-Y-(1+t)=0
\]
are coefficient-computable Puiseux surcomplex numbers. It does \emph{not}
provide a uniform labeling of these roots from arbitrary coefficient names.
Finding the appropriate finite ramification, identifying a root, and
certifying a separating finite jet are additional algorithmic tasks.

\subsection{Certificates rather than hidden zero tests}

The proof suggests a precise interface for an implicit root. A
\emph{simple-branch certificate} supplies a ramified coordinate $u$, an
integral polynomial $f$, an integer $s$, an integer $M>s$, and a polynomial
$a$ with computable real coefficients, with the obligations
\begin{equation}
 v_u(f(a))\geq M+s,\qquad v_u(f'(a))=s.
 \label{C-eq:branch-cert}
\end{equation}
These conditions imply that the transform \eqref{C-eq:Q-transform} reduces
to an affine polynomial with nonzero slope. The triangular algorithm then
constructs the unique root in $a+u^M\Rc[[u]]$. Thus a certificate can be
stated without first supplying an unknown full root.

The original proof shows that every simple formal root has some certificate
of this form. It does not show that all certificates can be found or
verified uniformly over arbitrary computable real coefficient programs.
A proof assistant may validate specific algebraic obligations, and a
restricted exact coefficient domain may decide them; neither fact makes
the general promise decidable.

\begin{proposition}[Unit-derivative lifting]
\label{C-prop:unit-hensel}
Let $F(u,Y)\in\Rc[[u]][Y]$ have uniformly computable coefficients, and let
$c\in\Rc$ satisfy
\[
 F(0,c)=0,\qquad F_Y(0,c)\ne0.
\]
Then there is a unique root $y\in\Rc[[u]]$ with constant term $c$, and its
coefficients are uniformly computable from these names and promises.
\end{proposition}
\begin{proof}
Set $y=c+\sum_{n\geq1}y_nu^n$. At degree $n\geq1$, the equation $F(u,y)=0$
is linear in the new unknown $y_n$, with coefficient $F_Y(0,c)$.
All other terms depend on finitely many previously computed coefficients.
Division by the fixed nonzero computable real $F_Y(0,c)$ therefore computes
$y_n$. Induction gives existence, uniqueness, and uniform computability.
This is the triangular mechanism of Lemma~\ref{C-lem:finite-jet} without the
preparatory change of scale.
\end{proof}

\begin{example}[Catalan branch]
The equation
\[
 y-t-y^2=0,\qquad [t^0]y=0,
\]
has unit derivative at $(t,y)=(0,0)$. Its coefficient recurrence is
\[
 y_1=1,\qquad y_n=\sum_{j=1}^{n-1}y_jy_{n-j}\quad(n\geq2),
\]
so
\[
 y=t+t^2+2t^3+5t^4+14t^5+42t^6+\cdots
   =\frac{1-\sqrt{1-4t}}2.
\]
The specified constant term chooses the branch. The other root has constant
term one and is not an alternative output for this branch certificate.
\end{example}

\begin{example}[A derivative that is not a unit]
Take
\[
 f(Y)=Y^2-u^2(1+u),\qquad a=u,
 \qquad s=1,\quad M=2.
\]
Then $v_u(f(a))=3=M+s$ and $v_u(f'(a))=1=s$. The certified transform is
\[
 u^{-3}f(u+u^2Z)=-1+2Z+uZ^2.
\]
Its reduction has the unique root $Z=1/2$. Lifting yields
\[
 y=u+\frac12u^2-\frac18u^3+\frac1{16}u^4
       -\frac5{128}u^5+\cdots=u\sqrt{1+u}.
\]
Thus the finite-jet method applies even when the original derivative has
positive valuation. The required condition is a sufficiently accurate seed
relative to that valuation, not an original unit derivative.
\end{example}


\paragraph{Coefficient-field generality.}
The finite-jet transformation is algebraic over any characteristic-zero
coefficient field. Its recursion uses uniform ring operations and the inverse
of only one nonzero slope. Thus if the coefficient field itself has only
extensional inverse closure, fix one name for that inverse as additional
finite seed data. This observation, emphasized in B, is exactly what the
iterated Puiseux proof needs; infinitely many independent inverse choices
would not establish a uniform coefficient sequence.

\section{Effective completion and the role of a modulus}
\label{sec:completion}

An ordinary valuation-Cauchy sequence eventually stabilizes every bounded
initial coefficient section. To turn this into an algorithm, one must know
\emph{when} stabilization has occurred.

\begin{definition}[Effective valuation-Cauchy sequence]
A uniformly numerically named sequence $(x_n)$ has a computable valuation
modulus if there is a total computable function $\mu:\Q\to\N$ such that
\begin{equation}\label{eq:modulus}
 n,m\ge\mu(B)\quad\Longrightarrow\quad
 \val(x_n-x_m)\ge B.
\end{equation}
This is an exact stabilization assertion for coefficients below $B$, not
a bound on their ordinary real approximation errors.
\end{definition}

\begin{theorem}[Effective sequential completeness]
\label{thm:completion}
From a uniformly named sequence in $\Lc$ and a computable valuation
modulus, one can compute its unique left-finite valuation limit as an
element of $\Lc$.
\end{theorem}
\begin{proof}
At exponent $q$, define the limiting coefficient by
\begin{equation}\label{eq:limcoef}
 a_q=\coef q{x_{\mu(q+1)}}.
\end{equation}
All later terms agree there, by the modulus at $q+1$. Thus this formula
supplies a joint Cauchy coefficient algorithm.

Let $n_0=\mu(0)$ and let $\ell_0$ be the supplied lower bound of $x_{n_0}$.
Set $\ell=\min(\ell_0,0)$. Every limiting coefficient below zero agrees
with $x_{n_0}$, so it vanishes below $\ell$. More explicitly, for any
$q<0$, choose an index later than both $\mu(0)$ and $\mu(q+1)$; the two
stabilization promises identify their coefficients at $q$.

For a threshold $B$, take $F_{x_{\mu(B)}}(B)$, filter it to $[\ell,B)$,
and use the resulting finite list as the limit cover. Every limiting
coefficient below $B$ agrees with $x_{\mu(B)}$: compare both to an index
later than $\mu(B)$ and $\mu(q+1)$. Hence the list is complete. The limit
has a numerical name, and for $n\ge\mu(B)$ it satisfies
$\val(x-x_n)\ge B$. Uniqueness follows by comparing every rational
coefficient, or by taking arbitrarily large $B$ in the ultrametric
inequality.
\end{proof}

\begin{theorem}[Effective completion characterization]
\label{thm:completioncharacter}
Every element of $\Lc$ is the limit of a uniformly computable sequence
of finite Puiseux polynomials with a computable valuation modulus.
Conversely, every such effective valuation limit of uniformly named
$\Pc$ elements belongs to $\Lc$.
\end{theorem}
\begin{proof}
Given $x\in\Lc$, take $p_n=J_nx$. Its computed exponent list is finite,
so the least common multiple of its denominators supplies a Puiseux grid.
The real coefficients have Cauchy names uniformly inherited from $x$.
Proposition~\ref{prop:jets} gives $\val(x-p_n)\ge n$, hence
$\val(p_n-p_m)\ge\min(n,m)$. Thus
$\mu(B)=\max(0,\lceil B\rceil)$ is a modulus. The converse follows by
translating Puiseux names to numerical names and applying
Theorem~\ref{thm:completion}.
\end{proof}

This is the precise completion analogy with computable real numbers in the
present theory. The finite approximants retain Cauchy-named real
coefficients, so it is a \emph{two-layer} completion, not the completion of
a discretely coded rational-monomial field in its valuation metric.
Rational constants cannot approximate an arbitrary irrational constant in
that metric, because every nonzero real difference has valuation zero.

\begin{proposition}[A computable Cauchy sequence without a computable limit]
\label{prop:nonmodulus}
There is a uniformly computable sequence of rational polynomials in $t$
which is valuation-Cauchy but whose left-finite limit has no numerical name.
It has no computable valuation modulus.
\end{proposition}
\begin{proof}
At stage $s$ put
\[
 z_s=\sum_{\substack{e\le s\\\text{machine }e\text{ halts by stage }s}}t^e.
\]
Each polynomial is computed by finitely many simulations. For a fixed
threshold $B$, only finitely many integers $e<B$ matter. Each one either
never halts or eventually has its coefficient permanently set to one.
Hence the sequence is valuation-Cauchy, with left-finite limit
\[
 z=\sum_{e\ge0}\chi_K(e)t^e.
\]
A numerical name for $z$ would compute its zero--one coefficients and
therefore decide $K$. If a computable modulus existed,
Theorem~\ref{thm:completion} would produce just such a name.
\end{proof}

Even a \emph{computable} modulus does not keep a sequence inside $\Pc$:
the finite partial sums of \eqref{eq:unboundeddenoms} have an obvious
computable modulus and converge to $u\in\Lc\setminus\Pc$. The support
cover, rather than a single denominator, is what survives effective
completion.


\subsection{A different limit theorem: coefficientwise real convergence}
\begin{theorem}[Uniform coefficientwise limits on a fixed lattice]
\label{thm:coefficient-limit}
Suppose $x_j=\sum_{n\ge N}a_{j,n}u^n$, $u=t^{1/d}$, share fixed $d,N$,
the real coefficients are uniformly computable in $(j,n)$, and a computable
function $M(n,k)$ satisfies
\[
 |a_{j,n}-a_{\ell,n}|\le 2^{-k}\qquad(j,\ell\ge M(n,k)).
\]
Then $a_n=\lim_j a_{j,n}$ is uniformly computable, and
$\sum_{n\ge N}a_nu^n$ has a uniformly constructed Puiseux name.
\end{theorem}
\begin{proof}
For $(n,k)$ choose $j=M(n,k+1)$ and approximate $a_{j,n}$ with error
$2^{-(k+1)}$. Its distance from the limit is at most $2^{-(k+1)}$ by real
completeness. The combined error is at most $2^{-k}$. Fixed metadata certify
the support of the result regardless of cancellation.
\end{proof}
This is not exact valuation stabilization: rational constants converging to
an irrational real never have positive valuation error. Without a modulus,
stagewise rational approximations to $r_K$ from
Proposition~\ref{prop:haltingreal} converge coefficientwise to a noncomputable
constant. Conversely, the exact-valuation example with increasingly many
halting coefficients has eventually \emph{equal} low coefficients but no
computable stabilization modulus. The two failures should not be identified.
Nor should $\Lc$ be called sequentially complete without qualification: the
computable-modulus theorem is strictly weaker than closure under every
classically Cauchy sequence of its elements.

\section{Sharp limits of uniform computation}
\label{sec:uniformlimits}

The field property and the effective algorithms already proved might suggest
that a missing normalization subroutine is merely a technical inconvenience.
The next results show an unavoidable obstruction.

\subsection{A uniformly computable real family encoding halting}

Fix a machine enumeration with halting stages numbered from one. Define
\begin{equation}\label{eq:ae}
 a_e=\begin{cases}
 2^{-s},&\text{machine }e\text{ first halts at stage }s,\\
 0,&\text{machine }e\text{ never halts}.
 \end{cases}
\end{equation}
This is a uniformly Cauchy-computable family of nonnegative reals. To
approximate $a_e$ to error $2^{-k}$, simulate for $k$ stages. If a halt has
occurred, return its value; otherwise return zero. A later halt contributes
less than $2^{-k}$.

\begin{theorem}[Valuation, inversion, and sign obstructions]
\label{thm:inversehard}
On raw numerical names, and already on raw Puiseux names:
\begin{enumerate}[label=\textup{(\roman*)},leftmargin=2em]
\item the valuation of a nonzero element is not uniformly computable;
\item reciprocation on nonzero inputs is not uniformly computable;
\item strict positivity and strict negativity are not semidecidable in
general, even on inputs promised nonzero;
\item equality is a $\Pi^0_1$-complete promise problem on valid indices,
although nonzeroness and disequality are semidecidable.
\end{enumerate}
\end{theorem}
\begin{proof}
For (i) use the always positive nonzero family
$x_e=a_e+t$. It has a fixed denominator-one name with support cover
contained in $\{0,1\}$, but
\[
 \val(x_e)=\begin{cases}0,&e\in K,\\1,&e\notin K.\end{cases}
\]
Thus computing the valuation would decide halting.

For (ii), the coefficient of the inverse at exponent $-1$ satisfies
\begin{equation}\label{eq:inversehalting}
 \coef{-1}{(a_e+t)^{-1}}=
 \begin{cases}0,&a_e>0,\\1,&a_e=0.\end{cases}
\end{equation}
When $a_e=0$ the inverse is $t^{-1}$. When $a_e>0$ its formal expansion
is $a_e^{-1}\sum_{n\ge0}(-t/a_e)^n$, which has no negative exponent.
A hypothetical inverse compiler followed by coefficient approximation to
error $1/8$ would decide $K$.

For (iii), put $y_e=a_e-t$. This is never zero, and
$y_e<0$ exactly when $e\notin K$. If negativity were semidecidable, the
complement of $K$ would be computably enumerable. Applying the same
argument to $-y_e$ excludes semidecidable positivity.

For (iv), equality to zero on the constants $a_e$ would decide whether
$e\notin K$. In contrast, to semidecide $x\ne0$, enumerate rational
exponents and precision levels in parallel. A coefficient is certified
nonzero whenever $|A(q,k)|>2^{-k}$. A nonzero coefficient eventually
produces such a certificate, while a zero series never does. Applying this to a difference semidecides
disequality ($x\ne y$). Equality is the absence of such a witness; the constants $a_e$
give $\Pi^0_1$-hardness. This classification concerns promised valid indices,
not the additional totality and consistency conditions for arbitrary programs.
\end{proof}

An important feature is that the inverse obstruction uses \emph{only}
nonzero, strictly positive inputs. Supplying a Boolean assertion
``nonzero'' or ``positive'' does not replace a leading-exponent certificate.

\begin{theorem}[A representation-independent incompatibility]
\label{thm:nogo}
Let $F$ be any represented subfield of a Hahn field containing the real
constants and a positive infinitesimal monomial $t$. Suppose that:
\begin{enumerate}[label=\textup{(\alph*)},leftmargin=2em]
\item ordinary real Cauchy names can be mapped uniformly to constant names;
\item $a\mapsto a+t$ is computable on those names;
\item the coefficient at $t^{-1}$ of an output element is uniformly
readable as a Cauchy real.
\end{enumerate}
Then no uniform reciprocal operation on all nonzero names can exist.
The conclusion already holds when the input values are the degree-one
polynomials $a+t$ with computable nonnegative real $a$.
\end{theorem}
\begin{proof}
Apply (a) and (b) to the uniform family \eqref{eq:ae}. An inverse operation
and (c), followed by \eqref{eq:inversehalting}, would decide halting.
\end{proof}

Enlarging the permitted support order types therefore cannot remove this
particular obstruction. A successful alternative representation must change
one of the capabilities: supply normalization data, weaken coefficient
access, use a more informative real representation, or treat some symbolic
operations as denotational rather than coefficient-computable.

\subsection{Exactly what one halting oracle buys}

\begin{theorem}[Normalization has halting degree on computable indices]
\label{thm:jump}
Given program indices for a valid computable numerical name, with nonzero
value promised, the halting oracle $0'$ suffices to compute its leading
exponent and to select an ordinary computable index for its inverse.
Conversely, a solver for either of these uniform tasks decides the halting
problem on computable inputs.
\end{theorem}
\begin{proof}
First search, without an oracle, for a certified nonzero coefficient.
The semidecision procedure in Theorem~\ref{thm:inversehard} terminates
because $x\ne0$. Let $q$ be its exponent. The finite cover $F(q+1)$
contains all nonzero coefficients at or below $q$ and only finitely many
other candidates. For each candidate $r$, the predicate $a_r\ne0$ is
computably enumerable uniformly from the input program index. A halting
oracle decides all these finitely many predicates. Taking the least
nonzero candidate yields the true leading exponent $\lambda$.

Now hard-code this rational $\lambda$ into the inverse compiler of
Theorem~\ref{thm:inverse}. The resulting inverse index needs no oracle
when it is subsequently run. The lower bounds are the explicit reductions
using $a_e+t$ in Theorem~\ref{thm:inversehard}.
\end{proof}

The qualification ``on computable indices'' is essential. A fixed oracle
$0'$ does not decide the zero status of arbitrary real coefficients supplied
by arbitrary noncomputable oracle names. In the relativized formulation,
the corresponding jump of the input oracle is used. We are not assigning
a more delicate Weihrauch degree to these operators here.


\subsection{Square roots, arbitrary root selection, and modulus}
\begin{theorem}[Root-selection obstructions]\label{thm:root-obstructions}
Neither positive square root on all positive bare names nor selection of any
root of every separable quadratic is uniformly computable in $\Pc$ or $\Lc$.
The first assertion holds even when a common positive lower bound $t$ is given.
\end{theorem}
\begin{proof}
For $x_e=a_e+t\ge t$, if $a_e=0$ its positive root is $t^{1/2}$.
If $a_e>0$ its positive root is
$\sqrt{a_e}\sum_{n\ge0}\binom{1/2}{n}(t/a_e)^n$, with integer exponents only.
The coefficient at $t^{1/2}$ is therefore one or zero according as $e$ fails
to halt or halts. Reading this coefficient decides halting.

For the stronger branch-independent assertion from B, consider
\[
 Q_e(Y)=Y^2-(a_e^2+t^2).
\]
This polynomial is separable even when $a_e=0$. In that case its two roots
are $\pm t$, with coefficient at $t$ equal to $\pm1$. If $a_e>0$, both roots
are $\pm a_e(1+t^2/a_e^2)^{1/2}$ and their coefficient at $t$ is zero.
An approximation with error less than $1/4$ to the coefficient of
\emph{any} selected root distinguishes the cases. The same argument applies
when the selector is allowed surcomplex outputs, since these quadratics
already split with exactly those two real roots.
\end{proof}

\begin{proposition}[Certified square root and the extended halting threshold]
\label{prop:root-jump}
For a positive numerical name with known true leading exponent $\lambda$,
positive square root is uniformly computable. On ordinary computable indices,
$0'$ is necessary and sufficient for equality, sign, valuation on nonzero
values, inverse on nonzero values, and positive square root on positive values.
After selecting finite leading data, the inverse and square-root outputs
can be ordinary computable programs.
\end{proposition}
\begin{proof}
Let $a=[t^\lambda]x>0$ and $h=a^{-1}t^{-\lambda}x-1$.
Then $h$ has positive support and a computable gap $\delta>0$, obtained
from its finite cover below one. Set
\[
 \sqrt{x}=\sqrt a\,t^{\lambda/2}
       \sum_{n\ge0}\binom{1/2}{n}h^n.
\]
Only finitely many powers can contribute below a given cutoff, since their
valuations are at least $n\delta$. Finite convolution and cover unions give
a valid numerical name. On a Puiseux grid this just doubles the denominator
if necessary and otherwise preserves the coefficient lattice. The formal
binomial identity proves that the result is the positive root.

A halting oracle decides the nonzeroness search for a whole series and for
any individual coefficient. Theorem~\ref{thm:jump} then finds the leading
exponent. Its nonzero real coefficient has a computable strict sign, and
the preceding construction gives the square root. The negative families
already displayed prove every lower bound. No upper bound for arbitrary
polynomial branch selection is asserted here.
\end{proof}

Leading data for two inputs do not automatically normalize their sum: leading
coefficients may cancel. Thus ring constructors return redundant names;
normalization is a separate task. A leading certificate is a mathematical
promise or a checked proof in a stated proof system, not a universally
decidable data-validation procedure.

\subsection{Exact coefficient domains change the uniformity boundary}

Suppose instead that coefficients lie in a computable ordered field $k$
with a decidable exact zero test, such as $\Q$ or the real algebraic
numbers with their usual exact descriptions. Keep the same effective
finite-section covers, but return exact coefficient codes.

\begin{proposition}[Exact coefficients permit leading-term search]
\label{prop:exactcoeffs}
For this exact-coefficient representation, the leading exponent and sign
of a nonzero series are computable uniformly. Reciprocation is uniformly
computable on nonzero names. Equality of arbitrary infinite series still
need not be decidable.
\end{proposition}
\begin{proof}
Inspect the covers $F(0),F(1),F(2),\ldots$ and their coefficients until a
nonzero coefficient is found. Every rational exponent is eventually below
one of these integer thresholds. At that stage, sort the finite cover and
choose its first nonzero entry; this is the true leading exponent because
all earlier possible nonzero positions are included. Exact order on $k$
gives its sign, and the geometric inverse construction is now uniform.

For equality, from a machine $e$ build the series whose coefficient at
integer $s\ge1$ is one exactly when $e$ first halts at stage $s$, and whose
other coefficients are zero. Coefficient access is a finite computation,
and the integer grid supplies covers. The series is zero exactly when the
machine never halts. Thus a uniform equality decider would decide halting.
\end{proof}

With $k$ equal to the real algebraic numbers the same Hensel and valuation
arguments prove an extensionally real-closed field. What is lost is the
uniform inclusion of arbitrary ordinary computable reals as constants,
not just algebraic constants. With $k=\Q$ the field is not real closed,
since a square root of the constant two cannot have rational coefficients.
These alternatives clarify which capabilities a symbolic system is buying
with an exact coefficient type.

\subsection{Countable does not mean effectively enumerable}

\begin{proposition}[A coefficient diagonal]
\label{prop:diagonal}
No uniformly computable sequence of valid numerical names enumerates all
of $\Lc$, even up to equality. The same is true for $\Pc$.
\end{proposition}
\begin{proof}
Let $(x_n)$ be such a sequence. Compute a rational $r_n$ within $1/8$ of
$\coef n{x_n}$, and choose
\[
 b_n=\begin{cases}1,&r_n\le1/2,\\0,&r_n>1/2.\end{cases}
\]
In the first case the true coefficient is at most $5/8$, and in the
second it is greater than $3/8$. Thus in either case it differs from $b_n$
by at least $3/8$. The series $y=\sum_{n\ge0}b_nt^n$ has a computable
Puiseux name, but $y\ne x_n$ for every $n$ by its coefficient at exponent
$n$. This proves both claims.
\end{proof}

The diagonal uses an approximation with a safety margin, not an equality
test on reals. It also shows why listing all programs and silently keeping
only the valid ones is not an executable enumeration strategy.

\section{Infinitesimal functional calculus and formal integration}
\label{sec:functions}

\subsection{Formal series need support bounds, not real convergence radii}

\begin{theorem}[Effective infinitesimal evaluation]
\label{thm:evaluation}
Let $f(Z)=\sum_{n\ge0}c_nZ^n$ have uniformly Cauchy-computable real
coefficients. If $h\in\Lc$ is numerically named with positive support,
then the Hahn evaluation
\[
 f(h)=\sum_{n\ge0}c_nh^n
\]
has a uniformly computable numerical name. If $h\in\Pc$, the result
belongs to $\Pc$ and its name can be computed on the same exponent grid.
\end{theorem}
\begin{proof}
Extract $\delta>0$ with $\val(h)\ge\delta$. A nonzero real $c_n$ has
valuation zero, so $\val(c_nh^n)\ge n\delta$. Below any threshold $B$,
only terms with $n\delta<B$ can contribute. Compute their finite covers
by the ring operations and take the union, after discarding negative
candidates as justified by the nonnegative support. At a queried exponent
$q$, choose $N$ with $(N+1)\delta>q$ and compute its coefficient in the
finite sum through $N$. This gives a total joint coefficient algorithm and
a lower bound zero. For a Puiseux input, finite sums, powers, and the
stabilized coefficients remain on its fixed grid.
\end{proof}

In particular $\exp(h)$, $\log(1+h)$, $\sin(h)$, $\cos(h)$, and
$(1+h)^a$ are computable whenever their formal coefficients are given by
the usual computable real formulas. The exponent $a$ in the binomial
series may itself be a Cauchy-computable real:
\[
 (1+h)^a=\sum_{n\ge0}\binom an h^n,
 \qquad \binom an=\frac{a(a-1)\cdots(a-n+1)}{n!}.
\]
For each coefficient, identities between these formal series reduce to
finite formal coefficient identities. They do not require convergence in
the full surreal order topology.

For a finite $x=r+h\in\Lc$ with $h$ infinitesimal and $r=\st(x)$,
one can define and compute
\[
 \exp_{\mathrm{fin}}(x)=e^r\exp(h).
\]
If $x$ has positive real residue $r$, then
$\log_{\mathrm{unit}}(x)=\log r+\log(1+(x-r)/r)$ is likewise computable
under that promise. These formulas are local, coefficientwise extensions
of ordinary real functions. No uniform computation of the unrestricted
surreal exponential, logarithm, or trigonometry at arbitrary infinite
arguments is asserted. Broader exponential and differential surreal
subfields involve additional closure questions and representations;
see, for example, \cite{bournez}.

\begin{example}[A ramified square root with exact infinitesimal terms]
Hensel lifting, or the formal binomial identity, gives
\[
 \sqrt{\omega^2+\omega}
 =\omega+\frac12-\frac1{8\omega}
  +\frac1{16\omega^2}-\frac5{128\omega^3}+\cdots.
\]
The positive branch is fixed by writing the result as
$\omega(1+t)^{1/2}$. For any exponent threshold, only finitely many
binomial coefficients are needed, and the remainder has a certified
valuation bound. The omitted terms are not an unspecified real error.
\end{example}

\subsection{A computable derivation and an exact obstruction to primitives}

Define a formal $t$-derivation by
\begin{equation}\label{eq:derivation}
 D_t\!\left(\sum_qa_qt^q\right)=\sum_q q a_qt^{q-1}.
\end{equation}
The coefficient field is held constant. This convention is important:
$D_t(\omega)=-\omega^2$, whereas the usual normalized
Berarducci--Mantova derivation has $D(\omega)=1$. We are not silently
identifying these two derivations or claiming a new construction of the
latter.

\begin{theorem}[Effective differentiation and the formal residue obstruction]
\label{thm:integration}
The map $D_t:\Lc\to\Lc$ is a uniformly computable derivation, with kernel
$\Rc$. Define
\[
 \res_t(x)=\coef{-1}x.
\]
There is a uniformly computable linear map $I_0:\Lc\to\Lc$ satisfying
\begin{equation}\label{eq:primitiveidentity}
 D_t I_0(x)=x-\res_t(x)t^{-1},\qquad
 \coef0{I_0(x)}=0.
\end{equation}
Consequently $x$ has a primitive in $\Lc$ exactly when $\res_t(x)=0$;
under that promise $I_0(x)$ is the unique primitive with zero constant
coefficient. The same statements hold in $\Pc$.
\end{theorem}
\begin{proof}
The coefficient at $r$ of $D_tx$ is $(r+1)a_{r+1}$. Shift a cover for $x$
below $B+1$ down by one and use lower bound $\ell_x-1$; optional deletion
of the exponent arising from $q=0$ is justified exactly. This is a numerical
name. For each coefficient of a product, the convolution sum is finite.
The identity $r+s$ for the derivative multiplier at a product term splits
as the sum of the two derivative multipliers, proving the Leibniz rule.
Characteristic zero implies that every coefficient at a nonzero input
exponent must vanish in the kernel. Thus the kernel consists exactly of
real constants in $\Lc$, namely $\Rc$.

Set
\begin{equation}\label{eq:I0}
 I_0(x)=\sum_{q\ne-1}\frac{a_q}{q+1}t^{q+1}.
\end{equation}
The coefficient at exponent $r\ne0$ is $a_{r-1}/r$, and the coefficient
at zero is set to zero. Shifting $F_x(B-1)$ by one and deleting zero
provides a finite cover, with lower bound $\ell_x+1$. Thus the construction
is uniform. Direct coefficient calculation gives
\eqref{eq:primitiveidentity}.

For every $y$, the coefficient at $-1$ of $D_ty$ is zero, since it could
only come from the exponent-zero term, multiplied by zero. This proves
necessity of the residue condition; \eqref{eq:primitiveidentity} proves
sufficiency and the kernel statement proves uniqueness. Rational shifting
preserves a fixed Puiseux grid, proving the final assertion.
\end{proof}

The condition $\res_t(x)=0$ is not generally decidable from a raw numerical
name, but the operator $I_0$ itself is total and computable. Thus a useful
interface can always return a primitive-with-obstruction pair
$(I_0(x),\res_t(x))$ without deciding whether the obstruction vanishes.
For $x=t^{-1}$ the obstruction is one; a primitive would require a
logarithmic extension rather than another left-finite rational-power
series.


\subsection{The full residue obstruction and a differential limitation}
\begin{corollary}[Kernel, image, and quotient of the formal derivative]
\label{cor:differential-quotient}
For $F=\Pc$ or $F=\Lc$, let $D=D_t$ and let $I_0$ be the zero-constant
primitive operator omitting the exponent $-1$. Then
\[
 I_0Dx=x-[t^0]x,\qquad DI_0x=x-\res_t(x)t^{-1},\qquad
 \ker D=\Rc,
\]
where $\res_t(x)=[t^{-1}]x$. Moreover,
\[
 D(F)=\ker(\res_t),\qquad F/D(F)\cong\Rc
\]
as $\Rc$-vector spaces. Every primitive of an integrable $x$ is $I_0x+c$,
$c\in\Rc$. Deciding whether a primitive exists is not computable on raw names.
\end{corollary}
\begin{proof}
The coefficient identities follow from $D(t^q)=qt^{q-1}$ and division by
$q+1$ away from $q=-1$. They prove both the kernel and image descriptions;
residue is onto since $ct^{-1}$ has residue $c$. Two primitives differ by
a member of the kernel. Finally $ct^{-1}$ is integrable exactly when the
computable real $c$ is zero, giving the decision obstruction.
\end{proof}
Here $D(F)$ denotes the additive image, not an ideal of fields; the quotient
is a vector-space quotient. A logarithm symbol could supply a primitive of
$t^{-1}$ in a larger differential field, but does not make that primitive
an element of the present coefficient field.

\begin{proposition}[A simple differential equation outside the workspace]
\label{prop:ode-obstruction}
The equation $Dy=-t^{-2}y$ has no nonzero solution in $\Pc$ or $\Lc$.
\end{proposition}
\begin{proof}
For $y\ne0$, termwise differentiation gives
$\val(Dy)\ge\val(y)-1$, allowing infinity when $Dy=0$.
The proposed right side has valuation $\val(y)-2$, a contradiction.
\end{proof}
Thus real closedness and effective infinitesimal calculus do not imply
closure under arbitrary differential equations. In a suitable larger
exponential differential field, $\exp(t^{-1})$ is the expected kind of
solution; it is not licensed here by the positive-support substitution rule.


\section{Surcomplex computability}
\label{sec:surcomplex}

Define the numerically computable surcomplex field by
\[
 \Lc^{\mathbb C}=\Lc[i]=\{x+iy:x,y\in\Lc\}.
\]
A name is a pair of numerical names. Equivalently, one can use uniformly
Cauchy-computable complex coefficients and the union of the two real support
covers. The complex valuation is $\val(x+iy)=\min(\val(x),\val(y))$,
with the usual zero convention.

\begin{proposition}[Effective complexification]
\label{prop:complexification}
The two naming conventions are uniformly interconvertible. Addition,
multiplication, conjugation, real and imaginary parts, and coefficient
extraction are computable. The field is algebraically closed and has
residue field $\Cc=\Rc[i]$. Reciprocation exists extensionally but remains
nonuniform on unrestricted nonzero names.
\end{proposition}
\begin{proof}
Convert between paired real oracles and complex coefficient oracles
componentwise. Support covers are united in one direction and reused in
the other, since a complex coefficient is zero exactly when both its real
and imaginary parts are zero. The formulas
\[
 (x+iy)(u+iv)=(xu-yv)+i(xv+yu),\qquad
 \overline{x+iy}=x-iy
\]
use only the uniform real ring operations. Algebraic closedness is
Theorem~\ref{thm:Lcrealclosed}. Coefficient zero gives the complex residue
homomorphism with residue field $\Cc$.
For $z\ne0$, $z^{-1}=\bar z/(x^2+y^2)$ belongs to the field by
extensional real field closure. A uniform inverse algorithm would
restrict to the real input family $a_e+t$ and contradict
Theorem~\ref{thm:nogo}.
\end{proof}

A supplied true leading exponent and its nonzero complex coefficient
restore the geometric inverse algorithm: ordinary complex reciprocation
is computable with the nonzero promise. Complex simple residue roots lift
by precisely the Newton proof in Theorem~\ref{thm:hensel}, since the
complex leading coefficient of each derivative remains the same nonzero
residue. Formal evaluation and $D_t,I_0,\res_t$ extend componentwise.

This gives a substantial exact surcomplex algebra, including algebraic
solutions, infinitesimal holomorphic power series, and residue-controlled
formal primitives. It is not a theorem that arbitrary contours or
analytic functions on the whole surcomplex plane have effective
representations. Such claims would require specifying function names,
support dependencies, and contour conventions separately.


\begin{proposition}[Squared modulus is uniform; modulus is not]
\label{prop:modulus}
On pair names in $\Pc[i]$ and $\Lc[i]$, $|z|^2$ is uniformly computable.
Every $|z|$ belongs to the real base field, but the map $z\mapsto|z|$ is not
uniformly computable, even on nonzero real inputs.
\end{proposition}
\begin{proof}
For $z=x+iy$, $|z|^2=x^2+y^2$ uses ring operations. Real closedness supplies
its nonnegative square root. For $z_e=t-a_e$, the coefficient at $t$ in
$|z_e|$ is $1$ when $a_e=0$, and $-1$ when $a_e>0$, since then
$|z_e|=a_e-t$. Uniformly reading that output coefficient would decide $K$.
\end{proof}
Certified nonzero leading data permit inverse and modulus constructions, but
adjoining $i$ does not restore a missing normalization algorithm. Formal
complex substitution is componentwise and support-controlled; no contour,
global analytic continuation, or unrestricted global exponential theorem
follows from algebraic closedness alone.

\section{Beyond one scale: an iterated numerical atlas}
\label{sec:higher}

The rational-exponent core is useful but omits $\omega^\omega$ and many
other structurally simple numbers. It can be extended without abandoning
ordinary computable real constants. The following construction makes that
claim precise rather than treating a higher-rank Hahn field as automatically
effective.

\subsection{An effective left-finite extension of a represented field}

Suppose $k$ is a countable ordered field all of whose elements have
computable names, with uniform named ring operations and extensional
inverse closure. Define
$\operatorname{ELF}(k;u)$ by series
\[
 X=\sum_{q\in\Q} a_q u^q,
\]
where a program, given $q$, returns a $k$-name for $a_q$, together with a
rational lower bound and computable finite covers of the nonzero outer
support. At the bottom level, returning a real name means returning an
index for a Cauchy program. At higher levels it means returning an index
and finite data for a name at the preceding level. Correctness and totality
are again semantic promises.

Finite sums and products use exactly the cover formulas from
Section~\ref{sec:arithmetic}, with finite coefficient arithmetic in $k$.
For inversion of a particular nonzero $X$, hard-code its leading rational
exponent \emph{and a computable name for the inverse of its leading
coefficient in $k$}. The geometric construction then works uniformly with
those data. This extra coefficient-inverse datum cannot be omitted when
$k$ itself has nonuniform inversion.

\begin{theorem}[Iterated numerically conservative real-closed fields]
\label{thm:iterated}
Put $\mathcal L_0=\Rc$ and inductively
\[
 \mathcal L_d=\operatorname{ELF}(\mathcal L_{d-1};t_d),\qquad
 t_d=\omega^{-\omega^{d-1}}\quad(d\ge1),
\]
with $t_d$ treated as an outer infinitesimal relative to $\mathcal L_{d-1}$.
Then each $\mathcal L_d$ embeds as a countable real-closed subfield of $\No$;
$\mathcal L_1=\Lc$; and the constant embeddings give
\[
 \mathcal L_0\subset \mathcal L_1\subset \mathcal L_2\subset\cdots.
\]
Their union $\mathcal L_{<\omega}$ is a countable real-closed subfield with
$\mathcal L_{<\omega}\cap\R=\Rc$. It contains every
$\omega^{\omega^{d-1}}$ for positive integer $d$.
\end{theorem}
\begin{proof}
Proceed by induction on $d$. The coefficient expansion of an element of
$\mathcal L_d$ flattens to a real Hahn series whose exponents lie in
\[
 \Gamma_d=\Q+\Q\omega+\cdots+\Q\omega^{d-1}\subseteq\No.
\]
The order on $\Gamma_d$ is lexicographic with its highest-degree
coordinate most significant. Each outer support is well ordered, and for
each of its exponents the inner coefficient support is well ordered.
A nonempty subset of the flattened support has a least outer coordinate,
then a least inner exponent. Induction gives well-ordering of the flattened
support. Uniqueness of the coordinate decomposition in $\Gamma_d$ and
Hahn arithmetic show that flattening is an ordered-field embedding into
$\R((\omega^{-\Gamma_d}))\subseteq\No$.

For clarity, inner exponents need not be bounded by real numbers. What
matters is that every inner exponent is smaller in magnitude than a
positive rational multiple of $\omega^{d-1}$. Thus $t_d$ is indeed smaller
than every positive element of the previously constructed field, as the
outer-series interpretation requires.

The finite-cover algorithms just described give uniform ring operations
and extensional field closure. Countability follows by induction from
countably many program indices. For the outer valuation, the value group
is $\Q$ and the residue field is $\mathcal L_{d-1}$, and its valuation ring is
convex for the lexicographic order.

We verify henselianity with special attention to the possible nonuniformity
in the coefficient field. Let $P$ have integral outer coefficients and
let $c\in \mathcal L_{d-1}$ be a simple residue root. Its derivative residue
$a\in \mathcal L_{d-1}$ is nonzero. By the inductive field property, select once
and for all a computable $\mathcal L_{d-1}$-name for $a^{-1}$. Every Newton
derivative then has the same outer leading coefficient $a$. Hence every
outer unit inverse can be computed with the same selected name for
$a^{-1}$, and no fresh nonuniform coefficient choices are needed at later
iterations. The positive-gap and doubling proof of
Theorem~\ref{thm:hensel} produces uniformly named stabilized outer
coefficients and a finite-cover algorithm for the root. This proves
henselianity. The inductive real closedness of $\mathcal L_{d-1}$ and
Lemma~\ref{lem:valuationRC} give real closedness of $\mathcal L_d$.

The constant embedding at each stage is effective. Every finite set of
coefficients from the union lies in some $\mathcal L_d$, where the requisite
positive square roots and odd-degree roots exist; hence the union is real
closed. Uniqueness of flattened Hahn normal forms shows that a real
constant in $\mathcal L_d$ must already have been a real constant in $\mathcal L_{d-1}$.
Induction gives the real-slice assertion. Finally $t_d^{-1}$ supplies
$\omega^{\omega^{d-1}}$.
\end{proof}

This atlas goes beyond a single fixed rank while retaining numerical
conservativity. It still does not cover all surreal scales or all
structurally computable surreals. Its exponent groups are finite rational
polynomials in $\omega$; for example, $\omega^{\omega^\omega}$ is not
in this particular union, by normal-form uniqueness. Adding a new scale
is a mathematical enlargement of the workspace, not an implicit
coercion that changes the existing exponent order.

The same construction works for complexifications. Each $\mathcal L_d[i]$ and
their union are algebraically closed. No uniform reciprocal claim is
restored by passing to this union: ordinary real inputs and the coefficient
of the original $t_1^{-1}$ still give the obstruction of
Theorem~\ref{thm:nogo}.


\subsection{The parallel Puiseux hierarchy and the common oracle bound}
Define $\mathcal P_0=\Rc$ and, with the same scales $t_d$, let
$\mathcal P_d$ consist of lower-bounded outer Puiseux series whose coefficient
names in $\mathcal P_{d-1}$ are generated by one total ordinary program.
The outer denominator is fixed for each name, not for the entire level.
Thus $\mathcal P_1=\Pc$, whereas $\mathcal L_1=\Lc$.

\begin{theorem}[Two compatible finite-rank hierarchies]
\label{thm:parallel-towers}
For every $d\ge1$, $\mathcal P_d$ is a countable real-closed subfield of $\No$,
with real trace $\Rc$, and there is a uniform embedding of names
$\mathcal P_d\hookrightarrow\mathcal L_d$. This inclusion is proper.
Both hierarchies have natural value group
$\Gamma_d=\Q+\Q\omega+\cdots+\Q\omega^{d-1}$, or equivalently lexicographic
$\Q^d$. Their outer valuation instead has value group $\Q$ and residue the
preceding level. Their finite-level unions are real closed and preserve
exactly the real trace $\Rc$.
\end{theorem}
\begin{proof}
The field and embedding arguments use common outer denominators, finite
convolution, and the reciprocal recurrence, choosing once the inverse of
the leading lower-level coefficient. Every outer lattice supplies a finite
cover, and induction converts its inner names. To prove real closedness,
apply classical Newton--Puiseux over the preceding real-closed field and
the finite-jet argument of Section~\ref{C-sec:realclosed}. A fixed finite
jet and one named inverse slope give a uniform tail over that coefficient
field. Separability and the minimal-polynomial argument finish the induction.
Alternatively, the outer simple-Hensel argument stays on a common grid.

Flattening and zero-coordinate extraction prove the valuation and trace
statements exactly as in Theorem~\ref{thm:iterated}. Proper inclusion is
witnessed at level $d$ by
$\sum_{n\ge1}t_d^{n+1/(n+1)}$: it has effective left-finite outer covers but
no bounded outer denominator. Inner coefficients cannot absorb a nonzero
outer coordinate. Finally every polynomial over either union has its finitely
many coefficients in one finite level, where the required algebraic roots
already exist.
\end{proof}

\begin{proposition}[Finite rank does not require iterated jumps]
\label{prop:tower-jump}
At every finite level of either hierarchy, nonzeroness is semidecidable.
On valid ordinary computable indices, $0'$ is necessary and sufficient for
sign, leading exponent vector, inverse, and positive square root on their
appropriate nonzero or positive domains. Oracle-assisted inverse and root
construction can return ordinary programs.
\end{proposition}
\begin{proof}
Dovetail searches through rational coefficient addresses down the finite
nesting depth until a nonzero real leaf is detected. This halts precisely
for nonzero series, so $0'$ decides zero at every level. Find any nonzero
outer coefficient, use a finite cover to find the least such coefficient
(or scan the discrete Puiseux grid), then descend to its leading inner
coefficient with the same oracle. The inverse and binomial recursions use
only finitely selected lower-level inverse or root names and uniform ring
operations. The nesting depth decreases, so the recursion terminates; it
does not call the jump of $0'$. The rank-one reductions embedded as constants
or as expressions in $t_1$ give the lower bounds.
\end{proof}
The distinction between the outer valuation and the full natural valuation
prevents a misleading statement that every higher-rank residue field is
$\Rc$: only the full natural residue is $\Rc$. Finite rank enlarges the
scale geometry, not the ordinary computational power of a name.

\section{Oracle information and support geometry are independent}
\label{sec:oracle-hierarchy}
For an oracle $X\subseteq\N$, write $\R_X$ for the reals with $X$-computable
fast names. Define $\mathcal P_X$ and $\mathcal L_X$ by relativizing,
respectively, the Puiseux names and the coefficient-and-cover names to $X$.
The oracle subscript is an oracle set, not the finite level index $d$ used
in the preceding section. In particular $\mathcal P_{\varnothing}=\Pc$ and
$\mathcal L_{\varnothing}=\Lc$.

\begin{theorem}[Relativization and exact degree embeddings]
\label{thm:oracle-fields}
For every $X$, both fields are countable and real closed, have real trace
$\R_X$, value group $\Q$, and residue field $\R_X$.
For arbitrary oracles $X,Y$,
\[
 \mathcal P_X\subseteq\mathcal P_Y\iff X\le_TY,
 \qquad
 \mathcal L_X\subseteq\mathcal L_Y\iff X\le_TY.
\]
Thus equivalent Turing degrees give equal fields, and incomparable degrees
give incomparable fields within each family.
\end{theorem}
\begin{proof}
All finite arithmetic, cover, Hensel, and finite-jet algorithms relativize.
The ordinary-real algebraic proof also relativizes: finite isolating data
and an $X$-computable strict-inequality search compute any fixed algebraic
real root. The field and real-closure proofs therefore apply with $\R_X$.
Countability follows from countably many oracle program indices for fixed
$X$, and coefficient-zero extraction proves the real-trace assertion.

If $X\le_TY$, simulate each $X$ query using $Y$, including cover queries.
Conversely the constant
$r_X=\sum_{n\ge0}2\chi_X(n)4^{-(n+1)}$ belongs to both fields for $X$.
Inclusion in the corresponding field for $Y$ makes $r_X$ a $Y$-computable
real. The separated-digit argument from Proposition~\ref{prop:haltingreal}
then recovers $X$ from $Y$. This proves both equivalences.
\end{proof}
Consequently each family contains strict chains indexed by
$\varnothing,0',0'',\ldots$. The exact degree embedding was explicit for
Puiseux fields in C; the left-finite statement here is its direct
coefficient-and-cover analogue, with the proof given above.

\begin{proposition}[Unrestricted oracle unions]
\label{prop:oracle-unions}
As $X$ ranges over all subsets of $\N$, the union of the $\mathcal P_X$ is
$\PP(\R)$, the full classical real Puiseux field. The union of the
$\mathcal L_X$ is the full classical real Levi--Civita field.
\end{proposition}
\begin{proof}
For a fixed classical Puiseux series, encode a countable joint table of
fast rational names of its real coefficients into one oracle. For a fixed
left-finite series, additionally encode a complete finite list of its support
below each rational cutoff and choose one rational support lower bound.
These are countable data, so one oracle suffices. The reverse inclusions
are immediate from the respective support conditions.
\end{proof}
The oracle union does not remove its support restriction. Even with unlimited
oracle information, the Puiseux union misses the unbounded-denominator
example in $\Lc$. Likewise no rank-one oracle change inserts $\omega^\omega$.
Conversely enlarging supports cannot make a noncomputable real constant
computable under unchanged ordinary Cauchy observations. Similar degree
embeddings hold at every finite rank by the same real-trace argument.
Matching the real trace of the structural field by allowing hyperarithmetic
oracles does not prove equality of surreal fields or a name conversion.

\subsection{General Hahn workspaces require effective decomposition data}

There is another route to larger supports. Let $G$ be a computable ordered
abelian group with a specified ordered embedding into the surreal
exponent group. A coefficient oracle on $G$ can be supplemented by explicit
finite \emph{decomposition covers}, rather than by left-finite initial
sections. The next observation states exactly the extra data required for
one important operation.

\begin{proposition}[An effective Neumann interface]
\label{prop:neumanninterface}
Let $h=\sum_{g\in G}a_g t^g$ have well-ordered positive support and
uniformly computable real coefficients. Suppose a procedure, given $g$,
returns a finite set $W(g)$ of ordered finite tuples
$(g_1,\ldots,g_n)$ with sum $g$, containing every such tuple whose entries
all lie in $\supp(h)$. The empty tuple is included at $g=0$.
Then for every uniformly computable real coefficient sequence $(c_n)$,
the coefficients of $\sum_{n\ge0}c_nh^n$ are uniformly computable.
\end{proposition}
\begin{proof}
Neumann's support lemma gives the mathematical Hahn summability for
well-ordered positive support \cite{neumann}. Compute
\[
 \coef g{\left(\sum_{n\ge0}c_nh^n\right)}
   =\sum_{(g_1,\ldots,g_n)\in W(g)} c_n a_{g_1}\cdots a_{g_n}.
\]
Every actual nonzero contribution occurs in the list. Extra tuples have a
zero coefficient factor and contribute zero. The list is a set, so there
is no duplicate counting. The displayed expression is a finite real
Cauchy computation. The empty product is one.
\end{proof}

The classical lemma supplies existence and finiteness of actual
contributions. The algorithm $W$ supplies their effective completeness.
One must not conflate those assertions. In a full named-workspace design,
one must also supply whatever support certificates are required for
output names and for subsequent operations. This proposition proves
coefficient computation, not an automatic effective representation theorem
for every Hahn field.

Theorem~\ref{thm:badproduct} demonstrates that even an apparently tame
support can fail to supply the needed effective finite search information.
The iterated construction supplies it by nested finite covers; a general
Hahn implementation must expose its own equivalent mechanism.

\section{Integration with the Surreal repository}
\label{sec:repo}

\subsection{Inspected revision and scope}

The repository was inspected at revision
\begin{center}
\texttt{4896a2808ce30e01b1c86ae3ba2295a64762d246}.
\end{center}
The reading included the root README and the files
\begin{quote}\small
\path{Surreal/Foundations/SignSequenceField.lean}\\
\path{Surreal/HahnSeries/Evaluation.lean}.
\end{quote}
Repository structure information was also inspected \cite{repo}. This is a
scoped source review, not a complete audit or a new build of the repository.

At this revision the sign-sequence carrier already has an ordered-field
structure transported through its proved equivalence with numeric games.
The field source explicitly separates the remaining genetic product-cut,
real-closedness, and Hahn normal-form obligations. Thus it would be
incorrect to describe multiplication and inversion on the mathematical
carrier as wholly unimplemented. It would be equally incorrect to infer
from those field instances that the numerical algorithms of this article
have been implemented.

The Hahn evaluation source already requires positive order explicitly and
proves the substituted-term interpretation, a geometric identity, and an
exact finite remainder. Its evaluator is a mathematical, noncomputable
Lean definition using existing Hahn infrastructure. This is the right
semantic layer for our algorithms, but it does not supply computable
finite-section covers or ordinary Cauchy coefficient procedures. The
missing bridge is \emph{effective data with correctness proofs}, not a
second proof that a formal geometric sum exists.

\subsection{A proof-carrying architecture}

A practical design should keep three layers separate.

\paragraph{Semantic values.}
Use the existing Hahn field as the primary semantic target for the
rank-one development. Prove the properties of $\Lc$ first there. Transport
them into the sign-sequence surreal carrier only after the normal-form
embedding is available. A theorem about a named Hahn series must not be
advertised as a compiled theorem about the full surreal carrier before
that bridge is proved.

\paragraph{Executable names.}
Represent the rational lower bound, the coefficient program, and the
finite-cover program as explicit data. A validity proposition relates those
programs to a semantic series. The representation should not force a
canonical nonzero-support list, since that would hide a real zero test.
For an initial executable prototype, exact rational or algebraic
coefficients are easier; ordinary Cauchy coefficients must later retain
the nonuniformity distinctions proved here.

\paragraph{Operation certificates.}
Expose separate certificates for positive support, unit leading term,
leading exponent, a simple residue root, and a valuation-Cauchy modulus.
These certificates express different facts and cannot safely be collapsed
into a single ``nonzero'' flag. In particular, a positive-support proof can
justify deleting nonpositive zero candidates, and a unit proof can justify
inversion with leading exponent zero. Both are useful even when a general
normalization procedure is impossible.

The following is a data-contract sketch, not compiled Lean code:
\par\noindent\begin{minipage}{\linewidth}
\begin{lstlisting}
NumericalName:
    lower : Rational
    coefficient(q, precision) : Rational
    cover(bound) : finite set of Rational
    validity : error bounds, lower support bound,
               and completeness of every finite cover

LeadingCertificate(name):
    exponent : Rational
    coefficient_at_exponent_is_nonzero
    all_earlier_coefficients_are_zero

PositiveSupportCertificate(name):
    all_nonpositive_coefficients_are_zero

HenselInput:
    integral_polynomial_names
    residue_root_name
    residue_root_equation
    residue_derivative_is_nonzero
\end{lstlisting}
\end{minipage}\par

An output certificate may be proved without being found by a general
search. For example, differentiating a formal series proves that its
$t^{-1}$ coefficient is zero; multiplying a unit by another unit gives a
unit certificate; and the Hensel invariant supplies the derivative's
leading coefficient at every Newton stage. Proof generation along these
specific constructions is much more realistic than universal
normalization.

\subsection{A dependency order for formalization}

The first layer should formalize finite covers, truncation, and the two
finite formulas \eqref{eq:coeffmul}--\eqref{eq:covermul}. These are
finitary facts independent of real closedness. Next come the positive-gap
lemma, effective geometric series, and certified inverse, with a clear
separation between program construction and semantic correctness.
Effective stabilization and simple Hensel lifting then provide the main
analytic-algebraic engine. Real closedness can be proved through an
available valuation criterion or through an explicitly formalized version
of Lemma~\ref{lem:valuationRC}; it should not be added as an unproved
instance.

The negative results are valuable formalization targets as well. The
reduction using $a_e+t$ is much smaller than a full computable-analysis
library and guards against designing an impossible API. The support-product
counterexample tests a different assumption and should remain separate.
Structural validity and the hyperarithmetic classification belong to a
higher-recursion-theory track, rather than being prerequisites for the
rank-one numerical algorithms.

No new Lean files are claimed to compile in this package. The supplied
Python suites check finite exact identities and constructions only.
A source-level proof in this article, an executable test, and a kernel-checked
Lean theorem are three different kinds of evidence.

\section{Comparison, boundaries, and further research}
\label{sec:outlook}

\subsection{A capability comparison}

\begin{center}
\small
\begin{tabular}{@{}>{\raggedright\arraybackslash}p{0.25\textwidth}>{\raggedright\arraybackslash}p{0.33\textwidth}>{\raggedright\arraybackslash}p{0.34\textwidth}@{}}
\toprule
Feature & Structural field $\Sc$ & Numerical core $\Lc$ \\
\midrule
Input meaning & Hereditary effective cut construction & Cauchy coefficients and finite support covers \\
Real slice & Hyperarithmetic reals & Ordinary computable reals \\
Scales & All computable ordinal scales, and more values & Rational powers of $\omega$; enlarged by the atlas \\
Equality & Not a general decision procedure & Not decidable \\
Ring operations & Uniform cut compilers & Uniform finite-cover algorithms \\
Field closure & Established extensional theorem & Proved here; inverse names need normalization data uniformly \\
Root theory & Established real closedness & Effective simple Hensel lifts and extensional real closedness \\
Convergence & No nontrivial small net in full-surreal topology & Effective valuation limits with a modulus \\
\bottomrule
\end{tabular}
\end{center}

The entry for structural scales does not claim that every surreal of
computable birthday has an effective structural presentation. Likewise,
the atlas is not asserted to exhaust structural computability. A complete
effective translation between arbitrary structural names, ordinal sign
presentations, and the numerical names above is not proved here. None of
the main numerical theorems relies on such a translation.

\subsection{What is established, and what is not}

The main theorem has now been proved. Its countability and real slice are
Proposition~\ref{prop:realslice}; uniform ring arithmetic is
Theorem~\ref{thm:ring}; certified and extensional inversion are
Theorem~\ref{thm:inverse}; the precise obstruction is
Theorem~\ref{thm:inversehard}; Hensel lifting and real closedness are
Theorems~\ref{thm:hensel} and \ref{thm:Lcrealclosed}; Puiseux closure and
effective completion are Theorems~\ref{thm:puiseuxclosed} and
\ref{thm:completioncharacter}; and the functional and surcomplex claims
are Theorem~\ref{thm:evaluation}, Theorem~\ref{thm:integration}, and
Proposition~\ref{prop:complexification}. The iterated atlas is an additional
extension, proved in Theorem~\ref{thm:iterated}.

Several tempting stronger assertions were deliberately excluded. We did
not prove a general uniform algebraic root-isolation algorithm on raw
Cauchy--Hahn names. We did not define a universal exact symbolic zero test.
We did not infer effective finite decomposition from abstract Hahn
summability. We did not confuse a computable copy of a birthday with
computable rank access to its signs. We did not claim that valuations,
coefficient approximation, and full-surreal neighborhoods describe the
same topology. Finally, the current normal-form bridge and the new
computability theory have not been machine-checked as part of this work.

\subsection{Research directions with precise targets}

The following are directions left by this article, not claims that the
corresponding questions have never been studied elsewhere.

\paragraph{Uniform translation complexity.}
Determine the exact information needed to convert a numerical normal-form
name into an effective sign presentation, and conversely on a stated
subclass. The computable-copy example shows that rank access must be part
of the specification. The structural classification supplies a natural
upper-level language, but does not by itself yield a practical name
translator.

\paragraph{Sharper operator degrees.}
The normalization and inverse tasks have halting degree in the computable
index setting. Their precise represented-space or Weihrauch classification
under different cover encodings may distinguish one zero test from finite
parallel tests or more elaborate search. Any such refinement should
preserve the elementary lower-bound family $a+t$.

\paragraph{Support certificates beyond left-finiteness.}
Design reusable certified decomposition structures for more general
computable ordered groups, especially supports with genuine transfinite
order type rather than finite-rank nested left-finite structure. A useful
target is closure under multiplication and geometric inversion with
explicit finite dependency output for each requested coefficient.
Theorem~\ref{thm:badproduct} is a minimal adversarial test.

\paragraph{Effective algebraic branch data.}
Extend the certified algorithms already proved here to broader informative
root data: a specified Newton slope, a simple initial root, a ramification
bound, or a certified factorization pattern. The simple Hensel theorem and the finite-jet branch certificate are
already successful contracts. Repeated roots should not be handled
by assuming a zero decision for arbitrary coefficients.

\paragraph{Larger transcendental workspaces.}
Extend the numerical atlas by explicitly named logarithmic and exponential
scales, while proving that the real slice remains $\Rc$ and specifying
finite dependency rules for all included operations. A finite symbolic
expression and an effective coefficient expansion are different output
formats; a system can support both without asserting that they are uniformly
interconvertible.

\subsection{Conclusion}

A workable concept of computable surreal number is necessarily
representation-sensitive. Structural presentations capture hereditary
construction and admit hyperarithmetic real information. Numerical
Cauchy--Hahn names capture effective observation of coefficients and their
support dependencies, recover exactly the ordinary computable reals on
the real axis, and support a substantial real-closed theory.

The constructive side is not small: the numerical fields contain
infinitesimals, infinite numbers, arbitrary computable formal coefficient
families on controlled supports, algebraic roots, effective valuation
limits, and a hierarchy of higher scales. The obstructions are equally
precise. Finite mathematical dependency does not imply an effective list;
field closure does not imply a uniform inverse compiler; and a finite
program for a transfinite structure does not necessarily compute its real
value. Respecting those distinctions produces algorithms whose
mathematical contracts can actually be met.

\appendix
\section{Finite algorithms and reproducibility}
\label{app:checks}

\subsection{Rerun checks from all three reports}
All three original verification programs are retained unchanged, alongside
a new cross-model suite. They use Python standard-library exact arithmetic
without network access. Rerun results:
\begin{center}
\begin{tabular}{@{}lll@{}}
\toprule
Suite & Reported unit & Rerun result\\
\midrule
A & 4,967 exact comparisons & PASS\\
B & 18 named unit tests & PASS\\
C & 266 exact checks & PASS\\
Merged & 4,321 cross-model exact checks & PASS\\
\bottomrule
\end{tabular}
\end{center}
Unlike test units are not added together. The merged suite compares dense
and sparse convolution, redundant finite covers, branch transforms, residue
identities, and denominator examples. Synthetic halting fixtures do not
compute the halting set.

These finite checks do not prove the infinite theorems or validate arbitrary
coefficient programs. No new Lean verification or production surreal
computer-algebra implementation is claimed.

From the package root:
\begin{lstlisting}
python3 code/run_all.py
latexmk -pdf -interaction=nonstopmode -halt-on-error article.tex
\end{lstlisting}

\subsection{A precise finite-dependency recipe}

For reference, the positive-support evaluation algorithm can be implemented
by the following mathematical pseudocode. All polynomial coefficients
below are real \emph{names}, and \texttt{finite\_arithmetic} invokes the
correct real precision subroutine rather than replacing those coefficients
by exact floating-point numbers.
\begin{lstlisting}
coefficient_of_evaluation(f, h, q, precision, positive_support_proof):
    candidates = h.cover(1)
    delta = min({1} union {r in candidates : r > 0})
    choose N >= 0 with (N + 1) * delta > q
    form the named finite sum S = sum(n = 0..N, f[n] * h^n)
    return S.coefficient(q, precision)

cover_of_evaluation(f, h, bound, positive_support_proof):
    extract delta as above
    form the finite set of n >= 0 with n * delta < bound
    for these n, compute covers of f[n] * h^n below bound
    return their union, filtered to [0, bound)
\end{lstlisting}

For inversion, first normalize using the \emph{supplied} leading exponent
and its nonzero coefficient, then call this recipe with $f_n=1$. For Hensel
lifting, every derivative is a certified unit, so its leading exponent is
zero without a normalization search. These are explicit finite
computations per output request, even though a whole infinite series is
never printed.


\clearpage
\section{Reconciliation ledger and source crosswalk}
\label{app:reconcile}
This appendix records the mathematical disposition of the three manuscripts.
Locations below use their original TeX line numbers, not the repaginated PDF.
A has 2,235 lines, B has 1,916, and C has 2,221.  The three source manuscripts
are \emph{not} redistributed with this report; the line numbers and the
recorded hashes identify them, and the disposition of every result is set out
below, but the texts themselves are not included here.
The JSON manifest records full hashes; upload suffixes alone are not treated
as version history. A supplied the continuous left-finite development; shared
introductions, repeated proofs, bibliographies, and implementation advice were
consolidated rather than concatenated.

\subsection{Actual corrections and necessary qualifications}
\paragraph{Nonzero polynomial hypothesis.}
A, the corollary immediately after its real-closedness theorem (source lines
1083--1093), says every real surreal root of a polynomial over $\Lc$ is in
$\Lc$. Literally this fails for the zero polynomial. The merged statement
explicitly requires a \emph{nonzero} polynomial, as C already does. This is a
corrected hypothesis, not merely a notation change.

\paragraph{Two fields, not two answers about one field.}
A's $\Lc$ includes effective finite covers and variable rational denominators;
B--C's $\Pc$ has one bounded denominator per value. The proper inclusion and
completion theorem reconcile their scope. A's effective completeness is
never used to negate B--C's Puiseux incompleteness, and is not upgraded to
unqualified sequential completeness of $\Lc$.

\paragraph{A source question is answered by another source.}
C's first concluding question asks for effective supports beyond bounded
Puiseux denominators with algebraic closure and finite dependency data.
The finite-cover left-finite construction answers the real-closed version
and its complexification answers the algebraically closed version. The
remaining support research target is moved beyond that established case.

\paragraph{Closure is not a uniform algorithm.}
Every proof that chooses a leading exponent, multiplicity, finite jet, or
coefficient inverse is labeled as either finite nonuniform existence data
or an additional algorithmic input. No claim of raw-name field computability
is inferred from real closedness. C's positive-square-root lower bound and
B's any-root quadratic lower bound are retained separately.

\paragraph{Two different uses of one jump.}
B's unrestricted convolution algorithm returns an oracle-computable
coefficient function; it need not yield an ordinary computable product.
A--C's inverse and square-root normalization uses $0'$ to choose finite data,
after which an ordinary program suffices. The convolution search is phrased
with a c.e. contributor set, not an unjustified decidable zero test on real
coefficients. All such statements concern valid ordinary computable indices;
relative names require the corresponding relative jump.

\paragraph{Support and valuation conventions.}
All three use $t=\omega^{-1}$; increasing $t$ exponents correspond to
 decreasing $\omega$ exponents. Candidate covers may include zeros and need
not be nested. A finite cardinality bound is not a finite candidate list.
At higher rank, full natural valuation and outer valuation are distinguished.
The formal residue $[t^{-1}]x$ is not standard part $[t^0]x$.

\paragraph{Limits, certificates, and historical status.}
Exact coefficient-name jets and rational coefficient approximations are
separate observations. Exact valuation stabilization and real coefficientwise
Cauchy moduli are separate limit hypotheses. Proof-bearing certificates do
not make arbitrary semantic validity decidable. The structural classification
is credited to established work, not re-announced as a new theorem; the
repository's remaining real-closedness task means a formalization obligation,
not an open question about the classical surreal field.

\subsection{Topic-by-topic disposition}
\small
\begin{longtable}{@{}p{.23\textwidth}p{.25\textwidth}p{.44\textwidth}@{}}
\toprule
Topic & Original locations & Merged disposition\\
\midrule
\endhead
Real names, size and topology & A 223--318; B opening conventions; C 242--401 &
One common fast-Cauchy baseline, intrinsic versus ambient topology, and a separate note on infinitary machines.\\
Structural field and halting real & A 319--501; B structural section; C 402--666 &
Common cut compiler and separated base-four encoding retained; classification imported with coding caveat.\\
Validity and birthdays & A 478--500; C 480--547 &
Full $\Pi^1_1$ upper bound retained from C; all-left reduction supplies hardness. Recursive-ordinal boundary made explicit.\\
Signed computable orders & A 502--556; B signed-order construction; C 667--727 &
C's theorem for every c.e. set retains exact Turing degree; the finite-block variants are consolidated.\\
Hahn convolution failures & A 557--643; B convolution and one-jump section; C 728--821 &
Both left-finite and singleton-fiber examples retained; B's dyadic template recorded; sharp $0'$ upper bound included.\\
Finite-cover numerical model & A 644--902; C convolution certificate &
A's definition and full arithmetic proofs retained. Complete finite candidate lists are the operational requirement.\\
Hensel and valuation proof & A 903--1094 &
Retained with the explicit defectless residue-characteristic-zero input; zero polynomial excluded in the root corollary.\\
Puiseux field and inclusion & A 1095--1192; B Puiseux field; C 822--1028 &
One definition, common-denominator ring formulas, certified reciprocal, exact trace, and proper inclusion.\\
Finite advice and algebraic branches & B simple Hensel and finite-advice normalization; C 1029--1236 &
One full finite-jet proof, independent of the valuation proof, with Laurent reduction and an explicit certificate.\\
Uniformity and zero tests & A 1295--1479; B root/sign obstructions; C 1237--1428 &
Strongest compatible lower bounds; full promised-index classification; exact-coefficient alternative retained.\\
Limits and completion & A 1193--1294; C 1556--1643 &
Effective completion in $\Lc$, failure in $\Pc$, and a separate coefficientwise real limit theorem.\\
Formal functions and equations & A 1480--1610; B differentiation/integration; C 1429--1555 &
Formal evaluation, implicit examples, residue quotient, and a differential non-solvability example; not global analysis.\\
Surcomplex & A 1611--1662; B complex section; C 1840--1891 &
Both complexifications; uniform squared modulus versus nonuniform modulus; branch-dependent operations qualified.\\
Finite rank & A 1663--1776; B iterated Puiseux construction; C 1644--1750 &
Parallel $\mathcal P_d$ and $\mathcal L_d$ towers with proper inclusions, correct residues, and a single reused $0'$.\\
Oracle hierarchies & C 1751--1839 &
Puiseux degree embedding retained and proved for left-finite names too; unrestricted oracle unions identified separately.\\
General supports and development plan & A 1777--2043; B support certificates and repository plan; C 1892--2044 &
General Neumann certificates retained; redundant-name API, branch contract, and remaining support questions unified.\\
Checks and proof provenance & A 2064--2235; B verification/bibliography; C 2045--2221 &
All original test suites rerun; counts kept separate; additional cross-model tests; original source hashes preserved.\\
\bottomrule
\end{longtable}
\normalsize

\subsection{Bibliographic and repository checks}
The structural definition and classification were checked directly in pages
4--6 of Hamkins's 2024 lecture notes and his 2025 announcement. Lurie's
publisher record confirms June 1998, volume 63(2), pages 337--371; its 2014
online posting is not a new publication date. The full Lurie and Harkleroad
papers were not independently read in this merge. The Newton--Puiseux input
was checked in the introductions of Rond and Kedlaya; the defect formula
was checked in Rzepka--Szewczyk, Section 2.1, equation (1). The classical
left-finite definition was checked in Restrepo Borrero--Shamseddine,
Section 1.1, version 2. None of those sources is cited as having proved the
particular representation-specific synthesis above.

The repository tree and selected files were read at commit
\begin{center}\texttt{4896a2808ce30e01b1c86ae3ba2295a64762d246}\end{center}
the same snapshot recorded by the reports. The root README, the sign-field source, and the
Hahn evaluation source agree on the distinction between mathematical
infrastructure and effective observation algorithms. No new commit,
repository write, or local Lean build is part of this deliverable.

\clearpage
\section{A proof-dependency ledger}

\begin{longtable}{@{}p{0.31\textwidth}p{0.61\textwidth}@{}}
\toprule
Result & Dependencies and status \\
\midrule
\endhead
Structural classification & External theorem; historical sources and modern authors' statements cited. Not reproved. \\
Structural halting real & Effective cut and ring construction; independent explicit reduction, with prior phenomenon acknowledged. \\
Computable-copy sign trap & Finite halting simulations, canonical real sign expansions, and a direct decoding argument. \\
Noncomputable Hahn product & Explicit rational supports and base-four uniqueness. Full proof provided. \\
Numerical ring and field & Definition of covers; finite Cauchy arithmetic; positive gap; geometric identity. Full proofs provided. \\
Effective Hensel lifting & Certified unit inversion; finite Taylor identity; explicit valuation doubling and coefficient stabilization. Full proof provided. \\
Numerical real closedness & Hensel lifting, the proved residue and value-group descriptions, and the explicitly cited Ostrowski degree formula. \\
Puiseux real closedness & Two independent proofs: the common-grid Hensel/valuation route and Newton--Puiseux plus the proved finite-jet uniformity step. \\
Effective completion & Exact stabilization modulus; finite covers from a selected sequence term. Full proof provided. \\
Uniform inverse obstruction & The uniformly Cauchy-computable real family $a_e$ and the $t^{-1}$ coefficient of $(a_e+t)^{-1}$. Full proof provided. \\
Higher-rank atlas & Nested finite covers, one fixed coefficient-inverse choice per Hensel problem, and induction on rank. Full proof provided. \\
Oracle hierarchies & Relativized algorithms and the explicit separated-digit real encoding; both support classes treated. \\
Formalization status & No new Lean verification claimed; finite Python tests are supplementary only. \\
\bottomrule
\end{longtable}

\begin{thebibliography}{99}
\raggedright

\bibitem{conway}
J.~H.~Conway, \emph{On Numbers and Games}, Academic Press, 1976;
second edition, A K Peters, 2001. Classical surreal cuts, arithmetic,
real closedness, and normal forms.

\bibitem{gonshor}
H.~Gonshor, \emph{An Introduction to the Theory of Surreal Numbers},
London Mathematical Society Lecture Note Series 110, Cambridge University
Press, 1986. Sign sequences, normal forms, and surreal exponentiation.

\bibitem{harkleroad}
L.~Harkleroad, ``Recursive surreal numbers,'' \emph{Notre Dame Journal of
Formal Logic} \textbf{31}(3) (1990), 337--345.
\href{https://doi.org/10.1305/ndjfl/1093635498}{doi:10.1305/ndjfl/1093635498}.
Historical bibliographic reference; the full text was not accessible in
this preparation.

\bibitem{lurie}
J.~Lurie, ``The effective content of surreal algebra,'' \emph{Journal of
Symbolic Logic} \textbf{63}(2) (1998), 337--371.
\href{https://doi.org/10.2307/2586835}{doi:10.2307/2586835}.
Published abstract and bibliographic metadata consulted; the 2014 online
posting date is not the article's publication year.

\bibitem{hamkins}
J.~D.~Hamkins, ``The computable surreal numbers,'' Notre Dame Logic
Seminar, 3 December 2024; reporting joint work with D.~Turetsky.
Author's announcement and six-page handwritten lecture notes:
\url{https://jdh.hamkins.org/the-computable-surreal-numbers-notre-dame-logic-seminar-december-2024/}.
Lecture notes:
\url{https://jdh.hamkins.org/wp-content/uploads/Lecture-Notes-Computable-surreal-numbers-ND-Dec-2024.pdf}.
The relevant pages, including the final classification statement, were
visually inspected.

\bibitem{turetsky}
D.~Turetsky and J.~D.~Hamkins, author contributions to ``What do we know
about the computable surreal numbers?'' \emph{MathOverflow}, 2024,
question 473211.
\url{https://mathoverflow.net/questions/473211/what-do-we-know-about-the-computable-surreal-numbers}.
Primary research discussion; used together with the later lecture
statement, not as a substitute for a refereed paper.

\bibitem{sacks}
G.~E.~Sacks, \emph{Higher Recursion Theory}, Perspectives in Mathematical
Logic, Springer, 1990. Background on hyperarithmeticity, recursive
well-orders, and well-founded recursive trees.

\bibitem{restrepo}
M.~Restrepo Borrero and K.~Shamseddine, ``On a Lebesgue-like integral over
the Levi-Civita field $\mathcal R$,'' arXiv:2506.13930, version 2, 2025.
\url{https://arxiv.org/abs/2506.13930}.
Only the introductory definition and classical context of the
Levi--Civita field in \S1.1 are used here; its integration theory is not
an input to our formal $t$-integration theorem.

\bibitem{rzepka}
A.~Rzepka and P.~Szewczyk, ``Defect extensions and a characterization of
tame fields,'' arXiv:2209.03308, version 1, 2022.
\url{https://arxiv.org/abs/2209.03308}.
See \S2.1, equation (1), for the degree formula and the
residue-characteristic-zero defectless case.

\bibitem{neumann}
B.~H.~Neumann, ``On ordered division rings,'' \emph{Transactions of the
American Mathematical Society} \textbf{66} (1949), 202--252.
Classical support finiteness for positive well-ordered series.

\bibitem{bournez}
O.~Bournez and Q.~Guilmant, ``Surreal fields stable under exponential,
logarithmic, derivative and anti-derivative functions,''
arXiv:2211.08396, 2022.
\url{https://arxiv.org/abs/2211.08396}.
Context for larger surreal function-closed subfields, not a source for the
numerical representation defined here.

\bibitem{repo}
V.~Reshetnikov, \emph{Surreal}, GitHub repository, inspected revision
\texttt{4896a2808ce30e01b1c86ae3ba2295a64762d246}.
\url{https://github.com/VladimirReshetnikov/Surreal/tree/4896a2808ce30e01b1c86ae3ba2295a64762d246}.
Selected sources: root \texttt{README.md},
\path{Surreal/Foundations/SignSequenceField.lean}, and
\path{Surreal/HahnSeries/Evaluation.lean}. Repository statements describe
that revision; no local Lean build was performed for this article.


\bibitem{hamkins2025}
J.~D.~Hamkins, ``The computable surreal numbers, Fudan University, July 2025,''
author's seminar announcement, posted 16 July 2025, for the talk of 23 July.
\url{https://jdh.hamkins.org/the-computable-surreal-numbers-fudan-university-july-2025/}.
Reports joint work with D.~Turetsky and explicitly acknowledges partial
rediscovery of earlier work of Lurie.

\bibitem{rond}
G.~Rond, ``About the algebraic closure of the field of power series in several
variables in characteristic zero,'' \emph{Journal of Singularities}
\textbf{16} (2017), 1--51; arXiv:1303.1921v4.
\url{https://arxiv.org/abs/1303.1921}.
Only the introductory one-variable Newton--Puiseux formulation is used.

\bibitem{kedlaya}
K.~S.~Kedlaya, ``The algebraic closure of the power series field in positive
characteristic,'' \emph{Proceedings of the American Mathematical Society}
\textbf{129} (2001), 3461--3470; arXiv:math/9810142v3.
\url{https://arxiv.org/abs/math/9810142}.
The characteristic-zero theorem is stated in its introduction; the positive
characteristic results are not used here.

\bibitem{galeotti}
L.~Galeotti and H.~Nobrega, ``Towards computable analysis on the generalised
real line,'' arXiv:1704.02884, 2017.
\url{https://arxiv.org/abs/1704.02884}.
Abstract consulted for the distinction between ordinary and infinitary machines.

\bibitem{faverjon}
C.~Faverjon and J.~Roques, ``Hahn series and Mahler equations: Algorithmic
aspects,'' \emph{Journal of the London Mathematical Society} \textbf{110}(1)
(2024), e12945. \url{https://doi.org/10.1112/jlms.12945};
\url{https://arxiv.org/abs/2412.04928}.
Abstract and metadata consulted for structured support-receptacle algorithms.

\bibitem{wiki}
Wikipedia contributors, ``Computable number,'' user-supplied orientation page,
consulted for the ordinary computable-real question, not used as a proof of
the representation-specific theorems.
\url{https://en.wikipedia.org/wiki/Computable_number}.

\end{thebibliography}
\end{document}
